% Chapter 10: Data Analysis and Result Interpretation

\section{Learning Objectives}

After completing this chapter, readers will be able to:
\begin{itemize}
    \item Implement comprehensive data preprocessing pipelines for A/B test data
    \item Detect and handle outliers using multiple statistical approaches
    \item Address missing data with appropriate imputation strategies
    \item Apply data transformation techniques to meet analysis assumptions
    \item Conduct thorough exploratory data analysis with professional visualizations
    \item Build automated statistical analysis pipelines with assumption checking
    \item Calculate and interpret effect sizes and confidence intervals
    \item Distinguish between statistical and practical significance
    \item Perform segment analysis to identify heterogeneous treatment effects
    \item Create executive summaries and technical reports for different audiences
    \item Implement reproducible analysis workflows using Pandas and NumPy
    \item Generate interactive visualizations with Plotly and Matplotlib
\end{itemize}

\section{Introduction: From Data to Decisions}

The quality of experimental conclusions depends not just on experimental design but on the rigor of data analysis and interpretation. Between collecting raw data and making business decisions lies a critical analytical journey requiring statistical expertise, programming proficiency, and business acumen.

This chapter presents a comprehensive framework for analyzing A/B test results, from initial data validation through final reporting. We emphasize:

\begin{itemize}
    \item \textbf{Reproducibility:} Analyses should be scripted, versioned, and automated
    \item \textbf{Robustness:} Results should withstand sensitivity analyses
    \item \textbf{Transparency:} Methods and limitations should be clearly documented
    \item \textbf{Context:} Statistical findings must connect to business impact
    \item \textbf{Communication:} Insights should be accessible to all stakeholders
\end{itemize}

Modern A/B testing analysis leverages powerful Python libraries:
\begin{itemize}
    \item \textbf{Pandas:} Data manipulation and transformation
    \item \textbf{NumPy:} Numerical computations
    \item \textbf{SciPy:} Statistical tests and distributions
    \item \textbf{Statsmodels:} Advanced statistical modeling
    \item \textbf{Matplotlib/Seaborn:} Static visualizations
    \item \textbf{Plotly:} Interactive dashboards
\end{itemize}

We'll build practical, production-ready analysis code throughout this chapter.

\section{Data Preprocessing}

Raw experimental data typically requires substantial preprocessing before analysis. This section covers systematic approaches to data cleaning, validation, and transformation.

\subsection{Data Quality Assessment}

Before any analysis, assess data quality systematically.

\lstset{language=Python, caption=Comprehensive Data Quality Framework}
\begin{lstlisting}
import pandas as pd
import numpy as np
from scipy import stats
from typing import Dict, List, Tuple
import warnings

class ABTestDataValidator:
    """
    Comprehensive data quality validation for A/B tests.

    Checks for:
    - Sample ratio mismatch
    - Assignment consistency
    - Temporal anomalies
    - Missing data patterns
    - Duplicate records
    - Invalid values
    """

    def __init__(self, data: pd.DataFrame,
                 variant_col: str = 'variant',
                 user_col: str = 'user_id',
                 timestamp_col: str = 'timestamp',
                 metric_cols: List[str] = None):
        """
        Initialize validator.

        Parameters:
        -----------
        data : DataFrame
            Experimental data
        variant_col : str
            Column indicating variant assignment
        user_col : str
            Column with unique user identifiers
        timestamp_col : str
            Column with event timestamps
        metric_cols : List[str]
            Columns containing metrics to analyze
        """
        self.data = data.copy()
        self.variant_col = variant_col
        self.user_col = user_col
        self.timestamp_col = timestamp_col
        self.metric_cols = metric_cols or []
        self.issues = []
        self.warnings = []

    def run_all_checks(self) -> Dict:
        """Run all validation checks."""
        results = {
            'sample_ratio_mismatch': self.check_sample_ratio_mismatch(),
            'assignment_consistency': self.check_assignment_consistency(),
            'temporal_stability': self.check_temporal_stability(),
            'missing_data': self.check_missing_data(),
            'duplicates': self.check_duplicates(),
            'value_ranges': self.check_value_ranges(),
            'novelty_effects': self.check_novelty_effects()
        }

        results['summary'] = {
            'total_issues': len(self.issues),
            'total_warnings': len(self.warnings),
            'passed': len(self.issues) == 0,
            'issues': self.issues,
            'warnings': self.warnings
        }

        return results

    def check_sample_ratio_mismatch(self,
                                   expected_ratios: Dict = None,
                                   alpha: float = 0.01) -> Dict:
        """
        Check for Sample Ratio Mismatch (SRM).

        Parameters:
        -----------
        expected_ratios : Dict
            Expected allocation ratios {variant: proportion}
            If None, assumes equal allocation
        alpha : float
            Significance level for SRM test

        Returns:
        --------
        Dict with SRM test results
        """
        variant_counts = self.data[self.variant_col].value_counts()
        n_variants = len(variant_counts)
        total_n = len(self.data)

        # Default to equal allocation
        if expected_ratios is None:
            expected_ratios = {v: 1/n_variants for v in variant_counts.index}

        # Calculate expected counts
        expected_counts = np.array([
            expected_ratios.get(v, 1/n_variants) * total_n
            for v in variant_counts.index
        ])

        # Chi-square test
        chi2_stat, p_value = stats.chisquare(
            variant_counts.values,
            expected_counts
        )

        # Check for SRM
        has_srm = p_value < alpha

        if has_srm:
            self.issues.append(
                f"Sample Ratio Mismatch detected (p={p_value:.6f})"
            )

        return {
            'chi2_statistic': chi2_stat,
            'p_value': p_value,
            'has_srm': has_srm,
            'observed_counts': dict(variant_counts),
            'expected_counts': dict(zip(variant_counts.index, expected_counts)),
            'percent_deviation': dict(
                (v, (variant_counts[v] - expected_counts[i]) /
                 expected_counts[i] * 100)
                for i, v in enumerate(variant_counts.index)
            )
        }

    def check_assignment_consistency(self) -> Dict:
        """
        Check if users have consistent variant assignments.

        Returns:
        --------
        Dict with consistency check results
        """
        # Count unique variants per user
        variants_per_user = self.data.groupby(
            self.user_col
        )[self.variant_col].nunique()

        inconsistent_users = variants_per_user[variants_per_user > 1]
        n_inconsistent = len(inconsistent_users)
        pct_inconsistent = n_inconsistent / len(variants_per_user) * 100

        if pct_inconsistent > 1:
            self.issues.append(
                f"{n_inconsistent} users ({pct_inconsistent:.2f}%) "
                f"with inconsistent assignments"
            )
        elif pct_inconsistent > 0:
            self.warnings.append(
                f"{n_inconsistent} users with inconsistent assignments"
            )

        return {
            'n_inconsistent_users': n_inconsistent,
            'pct_inconsistent': pct_inconsistent,
            'inconsistent_user_ids': list(inconsistent_users.index),
            'is_consistent': n_inconsistent == 0
        }

    def check_temporal_stability(self, window: str = 'D') -> Dict:
        """
        Check for temporal instability in metric distributions.

        Parameters:
        -----------
        window : str
            Pandas frequency string for grouping ('D', 'H', etc.)

        Returns:
        --------
        Dict with temporal stability results
        """
        self.data[self.timestamp_col] = pd.to_datetime(
            self.data[self.timestamp_col]
        )
        self.data['period'] = self.data[self.timestamp_col].dt.floor(window)

        temporal_stats = {}

        for metric in self.metric_cols:
            # Calculate metric by period and variant
            period_stats = self.data.groupby(
                ['period', self.variant_col]
            )[metric].agg(['mean', 'std', 'count']).reset_index()

            # Calculate coefficient of variation over time
            cv_by_variant = {}
            for variant in period_stats[self.variant_col].unique():
                variant_data = period_stats[
                    period_stats[self.variant_col] == variant
                ]
                mean_of_means = variant_data['mean'].mean()
                std_of_means = variant_data['mean'].std()
                cv = std_of_means / mean_of_means if mean_of_means != 0 else np.inf
                cv_by_variant[variant] = cv

            # Flag high variability
            max_cv = max(cv_by_variant.values())
            if max_cv > 0.5:
                self.warnings.append(
                    f"High temporal variability in {metric} (CV={max_cv:.2f})"
                )

            temporal_stats[metric] = {
                'cv_by_variant': cv_by_variant,
                'max_cv': max_cv,
                'period_data': period_stats
            }

        return temporal_stats

    def check_missing_data(self, threshold: float = 5.0) -> Dict:
        """
        Analyze missing data patterns.

        Parameters:
        -----------
        threshold : float
            Percentage threshold for flagging high missingness

        Returns:
        --------
        Dict with missing data analysis
        """
        missing_pct = (self.data.isnull().sum() / len(self.data) * 100)

        high_missing = missing_pct[missing_pct > threshold]

        if len(high_missing) > 0:
            self.warnings.append(
                f"High missing data in columns: {dict(high_missing)}"
            )

        # Check if missingness differs by variant
        missing_by_variant = {}
        for col in self.data.columns:
            if self.data[col].isnull().sum() > 0:
                missing_rates = self.data.groupby(
                    self.variant_col
                )[col].apply(lambda x: x.isnull().sum() / len(x) * 100)
                missing_by_variant[col] = dict(missing_rates)

        return {
            'missing_percentages': dict(missing_pct),
            'high_missing_columns': dict(high_missing),
            'missing_by_variant': missing_by_variant
        }

    def check_duplicates(self) -> Dict:
        """Check for duplicate records."""
        n_total = len(self.data)
        n_unique = len(self.data.drop_duplicates())
        n_duplicates = n_total - n_unique
        pct_duplicates = n_duplicates / n_total * 100

        if pct_duplicates > 0:
            self.warnings.append(
                f"{n_duplicates} duplicate records ({pct_duplicates:.2f}%)"
            )

        return {
            'n_duplicates': n_duplicates,
            'pct_duplicates': pct_duplicates,
            'has_duplicates': n_duplicates > 0
        }

    def check_value_ranges(self) -> Dict:
        """Check for values outside expected ranges."""
        range_issues = {}

        for metric in self.metric_cols:
            data_series = self.data[metric].dropna()

            # Basic range statistics
            min_val = data_series.min()
            max_val = data_series.max()
            mean_val = data_series.mean()
            std_val = data_series.std()

            # Check for negative values where inappropriate
            if 'revenue' in metric.lower() or 'time' in metric.lower():
                n_negative = (data_series < 0).sum()
                if n_negative > 0:
                    self.warnings.append(
                        f"{n_negative} negative values in {metric}"
                    )
                    range_issues[metric] = {'negative_values': n_negative}

            # Check for extreme outliers (> 5 std from mean)
            z_scores = np.abs(stats.zscore(data_series))
            n_extreme = (z_scores > 5).sum()
            if n_extreme > 0:
                self.warnings.append(
                    f"{n_extreme} extreme outliers in {metric} (>5 SD)"
                )
                if metric in range_issues:
                    range_issues[metric]['extreme_outliers'] = n_extreme
                else:
                    range_issues[metric] = {'extreme_outliers': n_extreme}

        return range_issues

    def check_novelty_effects(self, periods: int = 3) -> Dict:
        """
        Check for novelty effects (first-period differences).

        Parameters:
        -----------
        periods : int
            Number of time periods to examine

        Returns:
        --------
        Dict with novelty effect analysis
        """
        self.data[self.timestamp_col] = pd.to_datetime(
            self.data[self.timestamp_col]
        )
        self.data['period_rank'] = self.data.groupby(
            self.variant_col
        )[self.timestamp_col].rank(method='dense', pct=True)

        # Split into early and late periods
        early_cutoff = 1.0 / periods

        novelty_results = {}
        for metric in self.metric_cols:
            early_data = self.data[self.data['period_rank'] <= early_cutoff]
            late_data = self.data[self.data['period_rank'] > early_cutoff]

            early_means = early_data.groupby(self.variant_col)[metric].mean()
            late_means = late_data.groupby(self.variant_col)[metric].mean()

            # Compare early vs late
            relative_change = ((late_means - early_means) / early_means * 100)

            novelty_results[metric] = {
                'early_means': dict(early_means),
                'late_means': dict(late_means),
                'relative_change_pct': dict(relative_change)
            }

        return novelty_results

    def generate_report(self) -> str:
        """Generate human-readable validation report."""
        results = self.run_all_checks()

        report = []
        report.append("=" * 60)
        report.append("A/B TEST DATA VALIDATION REPORT")
        report.append("=" * 60)
        report.append("")

        # Summary
        summary = results['summary']
        status = "PASSED" if summary['passed'] else "FAILED"
        report.append(f"Overall Status: {status}")
        report.append(f"Total Issues: {summary['total_issues']}")
        report.append(f"Total Warnings: {summary['total_warnings']}")
        report.append("")

        # Issues
        if summary['issues']:
            report.append("CRITICAL ISSUES:")
            for issue in summary['issues']:
                report.append(f"  - {issue}")
            report.append("")

        # Warnings
        if summary['warnings']:
            report.append("WARNINGS:")
            for warning in summary['warnings']:
                report.append(f"  - {warning}")
            report.append("")

        # Detailed results
        report.append("DETAILED RESULTS:")
        report.append("")

        # SRM
        srm = results['sample_ratio_mismatch']
        report.append(f"Sample Ratio Mismatch Test: p-value = {srm['p_value']:.6f}")
        report.append(f"  Observed counts: {srm['observed_counts']}")
        report.append("")

        return "\n".join(report)

# Example usage
if __name__ == "__main__":
    # Load experimental data
    # data = pd.read_csv('experiment_data.csv')

    # Initialize validator
    # validator = ABTestDataValidator(
    #     data=data,
    #     variant_col='variant',
    #     user_col='user_id',
    #     timestamp_col='timestamp',
    #     metric_cols=['conversion', 'revenue', 'engagement_time']
    # )

    # Run validation
    # print(validator.generate_report())
    pass
\end{lstlisting}

\subsection{Outlier Detection and Treatment}

Outliers can substantially bias estimates, particularly for mean-based metrics. We need systematic approaches to detect and handle extreme values.

\subsubsection{Detection Methods}

\begin{enumerate}
    \item \textbf{Statistical methods:} Z-scores, IQR, modified Z-scores
    \item \textbf{Visual methods:} Box plots, scatter plots, histograms
    \item \textbf{Domain knowledge:} Business rules and plausibility checks
    \item \textbf{Robust statistics:} Median absolute deviation (MAD)
\end{enumerate}

\lstset{language=Python, caption=Comprehensive Outlier Detection}
\begin{lstlisting}
class OutlierDetector:
    """
    Multiple methods for outlier detection in A/B test data.
    """

    @staticmethod
    def detect_zscore(data: pd.Series, threshold: float = 3) -> pd.Series:
        """
        Detect outliers using Z-score method.

        Parameters:
        -----------
        data : Series
            Numerical data
        threshold : float
            Z-score threshold (typically 2.5-3)

        Returns:
        --------
        Boolean series indicating outliers
        """
        z_scores = np.abs(stats.zscore(data.dropna()))
        return pd.Series(z_scores > threshold, index=data.dropna().index)

    @staticmethod
    def detect_iqr(data: pd.Series, k: float = 1.5) -> pd.Series:
        """
        Detect outliers using IQR method.

        Parameters:
        -----------
        data : Series
            Numerical data
        k : float
            IQR multiplier (1.5 for outliers, 3 for extreme outliers)

        Returns:
        --------
        Boolean series indicating outliers
        """
        Q1 = data.quantile(0.25)
        Q3 = data.quantile(0.75)
        IQR = Q3 - Q1

        lower_bound = Q1 - k * IQR
        upper_bound = Q3 + k * IQR

        return (data < lower_bound) | (data > upper_bound)

    @staticmethod
    def detect_mad(data: pd.Series, threshold: float = 3.5) -> pd.Series:
        """
        Detect outliers using Median Absolute Deviation (MAD).
        More robust to outliers than Z-score.

        Parameters:
        -----------
        data : Series
            Numerical data
        threshold : float
            MAD threshold (typically 3.5)

        Returns:
        --------
        Boolean series indicating outliers
        """
        median = data.median()
        mad = np.median(np.abs(data - median))

        # Modified Z-score
        modified_z_scores = 0.6745 * (data - median) / mad

        return np.abs(modified_z_scores) > threshold

    @staticmethod
    def detect_isolation_forest(data: pd.DataFrame,
                               contamination: float = 0.1) -> pd.Series:
        """
        Detect outliers using Isolation Forest algorithm.

        Parameters:
        -----------
        data : DataFrame
            Numerical features
        contamination : float
            Expected proportion of outliers

        Returns:
        --------
        Boolean series indicating outliers
        """
        from sklearn.ensemble import IsolationForest

        clf = IsolationForest(contamination=contamination, random_state=42)
        predictions = clf.fit_predict(data)

        return pd.Series(predictions == -1, index=data.index)

    @staticmethod
    def detect_domain_specific(data: pd.Series,
                              metric_name: str,
                              business_rules: Dict = None) -> pd.Series:
        """
        Apply domain-specific outlier rules.

        Parameters:
        -----------
        data : Series
            Metric data
        metric_name : str
            Name of metric for rule selection
        business_rules : Dict
            Custom business rules {metric: {'min': val, 'max': val}}

        Returns:
        --------
        Boolean series indicating outliers
        """
        if business_rules and metric_name in business_rules:
            rules = business_rules[metric_name]
            min_val = rules.get('min', -np.inf)
            max_val = rules.get('max', np.inf)
            return (data < min_val) | (data > max_val)

        # Default rules for common metrics
        if 'revenue' in metric_name.lower():
            # Revenue should be non-negative and below some maximum
            max_revenue = data.quantile(0.999) * 2
            return (data < 0) | (data > max_revenue)

        if 'time' in metric_name.lower():
            # Time metrics should be non-negative
            max_time = data.quantile(0.999) * 2
            return (data < 0) | (data > max_time)

        if 'rate' in metric_name.lower() or 'ratio' in metric_name.lower():
            # Rates should be in [0, 1] or [0, 100]
            return (data < 0) | (data > 100)

        return pd.Series(False, index=data.index)

    @staticmethod
    def consensus_detection(data: pd.Series, methods: List[str] = None,
                          min_agreement: int = 2) -> pd.Series:
        """
        Combine multiple detection methods with voting.

        Parameters:
        -----------
        data : Series
            Numerical data
        methods : List[str]
            Methods to use: ['zscore', 'iqr', 'mad']
        min_agreement : int
            Minimum number of methods that must agree

        Returns:
        --------
        Boolean series indicating consensus outliers
        """
        if methods is None:
            methods = ['zscore', 'iqr', 'mad']

        detector = OutlierDetector()
        results = []

        if 'zscore' in methods:
            results.append(detector.detect_zscore(data))
        if 'iqr' in methods:
            results.append(detector.detect_iqr(data))
        if 'mad' in methods:
            results.append(detector.detect_mad(data))

        # Count agreements
        agreement_count = sum(results)

        return agreement_count >= min_agreement

# Example visualization
def plot_outlier_analysis(data: pd.Series, metric_name: str):
    """Create comprehensive outlier visualization."""
    import matplotlib.pyplot as plt
    import seaborn as sns

    fig, axes = plt.subplots(2, 2, figsize=(14, 10))

    detector = OutlierDetector()

    # 1. Box plot
    axes[0, 0].boxplot(data.dropna())
    axes[0, 0].set_ylabel(metric_name)
    axes[0, 0].set_title('Box Plot')
    axes[0, 0].grid(alpha=0.3)

    # 2. Histogram with outlier thresholds
    axes[0, 1].hist(data.dropna(), bins=50, alpha=0.7, edgecolor='black')

    # Add IQR bounds
    Q1 = data.quantile(0.25)
    Q3 = data.quantile(0.75)
    IQR = Q3 - Q1
    lower = Q1 - 1.5 * IQR
    upper = Q3 + 1.5 * IQR

    axes[0, 1].axvline(lower, color='red', linestyle='--',
                       label=f'IQR Lower: {lower:.2f}')
    axes[0, 1].axvline(upper, color='red', linestyle='--',
                       label=f'IQR Upper: {upper:.2f}')
    axes[0, 1].set_xlabel(metric_name)
    axes[0, 1].set_ylabel('Frequency')
    axes[0, 1].set_title('Distribution with IQR Bounds')
    axes[0, 1].legend()
    axes[0, 1].grid(alpha=0.3)

    # 3. Q-Q plot
    stats.probplot(data.dropna(), dist="norm", plot=axes[1, 0])
    axes[1, 0].set_title('Q-Q Plot (Normal Distribution)')
    axes[1, 0].grid(alpha=0.3)

    # 4. Detection method comparison
    methods = {
        'Z-score': detector.detect_zscore(data),
        'IQR': detector.detect_iqr(data),
        'MAD': detector.detect_mad(data)
    }

    method_counts = {k: v.sum() for k, v in methods.items()}

    axes[1, 1].bar(method_counts.keys(), method_counts.values())
    axes[1, 1].set_ylabel('Number of Outliers Detected')
    axes[1, 1].set_title('Outlier Detection Method Comparison')
    axes[1, 1].grid(alpha=0.3)

    plt.tight_layout()
    return fig
\end{lstlisting}

\subsubsection{Treatment Strategies}

Once outliers are detected, choose appropriate treatment:

\begin{enumerate}
    \item \textbf{Removal:} Delete outlier observations (use cautiously)
    \item \textbf{Winsorization:} Cap at percentiles (e.g., 1st and 99th)
    \item \textbf{Transformation:} Log, square root, or other transforms
    \item \textbf{Robust statistics:} Use median, trimmed mean, or Huber loss
    \item \textbf{Separate analysis:} Report both with and without outliers
\end{enumerate}

\lstset{language=Python, caption=Outlier Treatment Methods}
\begin{lstlisting}
class OutlierTreatment:
    """Methods for treating outliers in A/B test data."""

    @staticmethod
    def winsorize(data: pd.Series,
                  lower_percentile: float = 1,
                  upper_percentile: float = 99) -> Tuple[pd.Series, Dict]:
        """
        Winsorize data by capping at percentiles.

        Parameters:
        -----------
        data : Series
            Numerical data
        lower_percentile, upper_percentile : float
            Percentile thresholds

        Returns:
        --------
        Tuple of (winsorized data, metadata dict)
        """
        from scipy.stats.mstats import winsorize as scipy_winsorize

        lower_bound = np.percentile(data.dropna(), lower_percentile)
        upper_bound = np.percentile(data.dropna(), upper_percentile)

        winsorized = data.clip(lower=lower_bound, upper=upper_bound)

        n_capped_lower = (data < lower_bound).sum()
        n_capped_upper = (data > upper_bound).sum()
        n_capped_total = n_capped_lower + n_capped_upper

        metadata = {
            'lower_bound': lower_bound,
            'upper_bound': upper_bound,
            'n_capped_lower': n_capped_lower,
            'n_capped_upper': n_capped_upper,
            'n_capped_total': n_capped_total,
            'pct_capped': n_capped_total / len(data) * 100
        }

        return winsorized, metadata

    @staticmethod
    def log_transform(data: pd.Series, offset: float = 1) -> pd.Series:
        """
        Apply log transformation.

        Parameters:
        -----------
        data : Series
            Numerical data (must be positive)
        offset : float
            Value to add before log (handles zeros)

        Returns:
        --------
        Log-transformed series
        """
        return np.log(data + offset)

    @staticmethod
    def remove_outliers(data: pd.DataFrame,
                       metric_col: str,
                       method: str = 'iqr',
                       **kwargs) -> Tuple[pd.DataFrame, Dict]:
        """
        Remove outlier observations.

        Parameters:
        -----------
        data : DataFrame
            Experimental data
        metric_col : str
            Column to check for outliers
        method : str
            Detection method: 'zscore', 'iqr', 'mad'
        **kwargs : additional arguments for detection method

        Returns:
        --------
        Tuple of (cleaned DataFrame, metadata dict)
        """
        detector = OutlierDetector()

        if method == 'zscore':
            outliers = detector.detect_zscore(data[metric_col], **kwargs)
        elif method == 'iqr':
            outliers = detector.detect_iqr(data[metric_col], **kwargs)
        elif method == 'mad':
            outliers = detector.detect_mad(data[metric_col], **kwargs)
        else:
            raise ValueError(f"Unknown method: {method}")

        n_outliers = outliers.sum()
        pct_removed = n_outliers / len(data) * 100

        cleaned_data = data[~outliers]

        metadata = {
            'method': method,
            'n_outliers': n_outliers,
            'pct_removed': pct_removed,
            'original_n': len(data),
            'cleaned_n': len(cleaned_data)
        }

        return cleaned_data, metadata

    @staticmethod
    def trimmed_mean(data: pd.Series, trim_pct: float = 0.1) -> float:
        """
        Calculate trimmed mean (removes extremes from both tails).

        Parameters:
        -----------
        data : Series
            Numerical data
        trim_pct : float
            Percentage to trim from each tail (0.1 = 10%)

        Returns:
        --------
        Trimmed mean value
        """
        from scipy.stats import trim_mean
        return trim_mean(data.dropna(), trim_pct)

    @staticmethod
    def compare_treatments(data: pd.DataFrame,
                          metric_col: str,
                          variant_col: str = 'variant') -> pd.DataFrame:
        """
        Compare effect estimates under different treatments.

        Parameters:
        -----------
        data : DataFrame
            Experimental data
        metric_col : str
            Metric to analyze
        variant_col : str
            Variant column

        Returns:
        --------
        DataFrame comparing treatment methods
        """
        treatment = OutlierTreatment()

        results = []

        # Original data
        for variant in data[variant_col].unique():
            variant_data = data[data[variant_col] == variant][metric_col]
            results.append({
                'treatment': 'original',
                'variant': variant,
                'mean': variant_data.mean(),
                'median': variant_data.median(),
                'std': variant_data.std(),
                'n': len(variant_data)
            })

        # Winsorized
        data_wins = data.copy()
        data_wins[f'{metric_col}_wins'], _ = treatment.winsorize(
            data[metric_col]
        )
        for variant in data_wins[variant_col].unique():
            variant_data = data_wins[
                data_wins[variant_col] == variant
            ][f'{metric_col}_wins']
            results.append({
                'treatment': 'winsorized',
                'variant': variant,
                'mean': variant_data.mean(),
                'median': variant_data.median(),
                'std': variant_data.std(),
                'n': len(variant_data)
            })

        # Log-transformed
        if (data[metric_col] > 0).all():
            data_log = data.copy()
            data_log[f'{metric_col}_log'] = treatment.log_transform(
                data[metric_col]
            )
            for variant in data_log[variant_col].unique():
                variant_data = data_log[
                    data_log[variant_col] == variant
                ][f'{metric_col}_log']
                results.append({
                    'treatment': 'log_transform',
                    'variant': variant,
                    'mean': variant_data.mean(),
                    'median': variant_data.median(),
                    'std': variant_data.std(),
                    'n': len(variant_data)
                })

        # Outliers removed
        data_clean, _ = treatment.remove_outliers(
            data, metric_col, method='iqr'
        )
        for variant in data_clean[variant_col].unique():
            variant_data = data_clean[
                data_clean[variant_col] == variant
            ][metric_col]
            results.append({
                'treatment': 'outliers_removed',
                'variant': variant,
                'mean': variant_data.mean(),
                'median': variant_data.median(),
                'std': variant_data.std(),
                'n': len(variant_data)
            })

        return pd.DataFrame(results)
\end{lstlisting}

\subsection{Missing Data Handling}

Missing data is ubiquitous in A/B tests. The approach depends on the missingness mechanism:

\begin{itemize}
    \item \textbf{MCAR (Missing Completely At Random):} Missingness independent of any variable
    \item \textbf{MAR (Missing At Random):} Missingness depends on observed variables
    \item \textbf{MNAR (Missing Not At Random):} Missingness depends on unobserved variables
\end{itemize}

\lstset{language=Python, caption=Missing Data Analysis and Imputation}
\begin{lstlisting}
class MissingDataHandler:
    """
    Comprehensive missing data analysis and treatment.
    """

    @staticmethod
    def analyze_missingness(data: pd.DataFrame) -> pd.DataFrame:
        """
        Comprehensive missing data analysis.

        Returns:
        --------
        DataFrame with missingness statistics per column
        """
        missing_stats = []

        for col in data.columns:
            n_missing = data[col].isnull().sum()
            pct_missing = n_missing / len(data) * 100

            missing_stats.append({
                'column': col,
                'n_missing': n_missing,
                'pct_missing': pct_missing,
                'dtype': str(data[col].dtype)
            })

        return pd.DataFrame(missing_stats).sort_values(
            'pct_missing', ascending=False
        )

    @staticmethod
    def test_mcar(data: pd.DataFrame,
                  target_col: str,
                  predictor_cols: List[str] = None) -> Dict:
        """
        Test if data is Missing Completely At Random using Little's test.

        Parameters:
        -----------
        data : DataFrame
            Data with missing values
        target_col : str
            Column to test for MCAR
        predictor_cols : List[str]
            Columns to use as predictors

        Returns:
        --------
        Dict with test results
        """
        if predictor_cols is None:
            predictor_cols = [c for c in data.columns if c != target_col]

        # Create missingness indicator
        data_test = data[predictor_cols + [target_col]].copy()
        data_test['is_missing'] = data_test[target_col].isnull().astype(int)

        # Test if predictors differ between missing and non-missing
        results = {}
        for col in predictor_cols:
            if data_test[col].dtype in ['int64', 'float64']:
                # T-test for numerical predictors
                missing_group = data_test[
                    data_test['is_missing'] == 1
                ][col].dropna()
                present_group = data_test[
                    data_test['is_missing'] == 0
                ][col].dropna()

                if len(missing_group) > 0 and len(present_group) > 0:
                    t_stat, p_val = stats.ttest_ind(
                        missing_group, present_group
                    )
                    results[col] = {
                        'test': 't-test',
                        'statistic': t_stat,
                        'p_value': p_val,
                        'differs_significantly': p_val < 0.05
                    }

        # Overall assessment
        n_significant = sum(
            1 for r in results.values() if r.get('differs_significantly', False)
        )

        likely_mcar = n_significant == 0

        return {
            'predictor_tests': results,
            'n_significant_predictors': n_significant,
            'likely_mcar': likely_mcar
        }

    @staticmethod
    def simple_imputation(data: pd.DataFrame,
                         col: str,
                         method: str = 'mean',
                         variant_col: str = None) -> pd.Series:
        """
        Simple imputation methods.

        Parameters:
        -----------
        data : DataFrame
            Data with missing values
        col : str
            Column to impute
        method : str
            'mean', 'median', 'mode', or 'forward_fill'
        variant_col : str
            If provided, impute within variants

        Returns:
        --------
        Series with imputed values
        """
        if variant_col:
            # Impute within each variant
            result = data[col].copy()
            for variant in data[variant_col].unique():
                mask = data[variant_col] == variant
                variant_data = data.loc[mask, col]

                if method == 'mean':
                    fill_value = variant_data.mean()
                elif method == 'median':
                    fill_value = variant_data.median()
                elif method == 'mode':
                    fill_value = variant_data.mode()[0]
                else:
                    raise ValueError(f"Unknown method: {method}")

                result.loc[mask] = variant_data.fillna(fill_value)

            return result
        else:
            # Global imputation
            if method == 'mean':
                return data[col].fillna(data[col].mean())
            elif method == 'median':
                return data[col].fillna(data[col].median())
            elif method == 'mode':
                return data[col].fillna(data[col].mode()[0])
            elif method == 'forward_fill':
                return data[col].fillna(method='ffill')
            else:
                raise ValueError(f"Unknown method: {method}")

    @staticmethod
    def multiple_imputation(data: pd.DataFrame,
                           cols_to_impute: List[str],
                           n_imputations: int = 5) -> List[pd.DataFrame]:
        """
        Multiple imputation using MICE algorithm.

        Parameters:
        -----------
        data : DataFrame
            Data with missing values
        cols_to_impute : List[str]
            Columns to impute
        n_imputations : int
            Number of imputed datasets to create

        Returns:
        --------
        List of imputed DataFrames
        """
        from sklearn.experimental import enable_iterative_imputer
        from sklearn.impute import IterativeImputer

        imputer = IterativeImputer(
            max_iter=10,
            random_state=42,
            sample_posterior=True
        )

        imputed_datasets = []

        for i in range(n_imputations):
            data_imputed = data.copy()
            imputer.set_params(random_state=42 + i)

            # Impute numerical columns
            numerical_cols = data[cols_to_impute].select_dtypes(
                include=['float64', 'int64']
            ).columns

            if len(numerical_cols) > 0:
                imputed_values = imputer.fit_transform(
                    data[numerical_cols]
                )
                data_imputed[numerical_cols] = imputed_values

            imputed_datasets.append(data_imputed)

        return imputed_datasets

    @staticmethod
    def analyze_with_multiple_imputation(
        imputed_datasets: List[pd.DataFrame],
        variant_col: str,
        metric_col: str
    ) -> Dict:
        """
        Analyze with multiple imputation and pool results (Rubin's rules).

        Parameters:
        -----------
        imputed_datasets : List[DataFrame]
            List of imputed datasets
        variant_col : str
            Variant column
        metric_col : str
            Metric to analyze

        Returns:
        --------
        Dict with pooled results
        """
        estimates = []
        variances = []

        for data_imp in imputed_datasets:
            # Calculate treatment effect in this imputed dataset
            variants = data_imp[variant_col].unique()
            if len(variants) == 2:
                control = variants[0]
                treatment = variants[1]

                control_data = data_imp[
                    data_imp[variant_col] == control
                ][metric_col]
                treatment_data = data_imp[
                    data_imp[variant_col] == treatment
                ][metric_col]

                # Difference in means
                diff = treatment_data.mean() - control_data.mean()

                # Variance of difference
                var_diff = (
                    control_data.var() / len(control_data) +
                    treatment_data.var() / len(treatment_data)
                )

                estimates.append(diff)
                variances.append(var_diff)

        # Pool using Rubin's rules
        m = len(imputed_datasets)

        # Pooled estimate
        pooled_estimate = np.mean(estimates)

        # Within-imputation variance
        within_var = np.mean(variances)

        # Between-imputation variance
        between_var = np.var(estimates, ddof=1)

        # Total variance
        total_var = within_var + (1 + 1/m) * between_var

        # Standard error
        se = np.sqrt(total_var)

        # Degrees of freedom (Barnard-Rubin adjustment)
        df = (m - 1) * (1 + within_var / ((1 + 1/m) * between_var)) ** 2

        # Confidence interval
        t_crit = stats.t.ppf(0.975, df)
        ci_lower = pooled_estimate - t_crit * se
        ci_upper = pooled_estimate + t_crit * se

        return {
            'pooled_estimate': pooled_estimate,
            'standard_error': se,
            'ci_lower': ci_lower,
            'ci_upper': ci_upper,
            'within_variance': within_var,
            'between_variance': between_var,
            'total_variance': total_var,
            'degrees_of_freedom': df,
            'individual_estimates': estimates
        }
\end{lstlisting}

\subsection{Data Transformation Techniques}

Transformations can improve adherence to statistical assumptions (normality, homoscedasticity).

\begin{enumerate}
    \item \textbf{Log transformation:} For right-skewed data ($y' = \log(y)$)
    \item \textbf{Square root:} For count data ($y' = \sqrt{y}$)
    \item \textbf{Box-Cox:} General power transformation
    \item \textbf{Logit:} For proportions ($y' = \log(p/(1-p))$)
    \item \textbf{Standardization:} Zero mean, unit variance ($y' = (y - \mu)/\sigma$)
\end{enumerate}

\lstset{language=Python, caption=Data Transformation Functions}
\begin{lstlisting}
class DataTransformer:
    """Data transformation methods for A/B test analysis."""

    @staticmethod
    def log_transform(data: pd.Series, offset: float = 1) -> pd.Series:
        """Log transformation with offset for zeros."""
        return np.log(data + offset)

    @staticmethod
    def sqrt_transform(data: pd.Series) -> pd.Series:
        """Square root transformation for count data."""
        return np.sqrt(data)

    @staticmethod
    def boxcox_transform(data: pd.Series) -> Tuple[pd.Series, float]:
        """
        Box-Cox transformation.

        Returns:
        --------
        Tuple of (transformed data, lambda parameter)
        """
        from scipy.stats import boxcox

        # Box-Cox requires positive data
        if (data <= 0).any():
            offset = abs(data.min()) + 1
            data_positive = data + offset
        else:
            data_positive = data

        transformed, lambda_param = boxcox(data_positive)

        return pd.Series(transformed, index=data.index), lambda_param

    @staticmethod
    def standardize(data: pd.Series) -> pd.Series:
        """Z-score standardization."""
        return (data - data.mean()) / data.std()

    @staticmethod
    def recommend_transformation(data: pd.Series) -> str:
        """
        Recommend transformation based on data characteristics.

        Returns:
        --------
        str : Recommended transformation
        """
        # Check skewness
        skew = stats.skew(data.dropna())

        # Check if data has zeros
        has_zeros = (data == 0).any()

        # Check if all positive
        all_positive = (data > 0).all()

        # Check range
        data_range = data.max() - data.min()
        mean_val = data.mean()

        # Decision tree for recommendation
        if abs(skew) < 0.5:
            return "none (data is approximately symmetric)"

        if skew > 1 and all_positive:
            return "log (right-skewed, positive data)"

        if skew > 1 and has_zeros:
            return "log_with_offset (right-skewed with zeros)"

        if skew < -1:
            return "reflect_and_log (left-skewed)"

        if 0.5 <= abs(skew) <= 1:
            return "sqrt (moderately skewed)"

        return "boxcox (general power transformation)"

    @staticmethod
    def assess_normality(data: pd.Series) -> Dict:
        """
        Assess normality using multiple tests.

        Returns:
        --------
        Dict with normality test results
        """
        # Shapiro-Wilk test (good for n < 5000)
        if len(data) < 5000:
            shapiro_stat, shapiro_p = stats.shapiro(data.dropna())
        else:
            shapiro_stat, shapiro_p = None, None

        # Anderson-Darling test
        anderson_result = stats.anderson(data.dropna())

        # Kolmogorov-Smirnov test
        ks_stat, ks_p = stats.kstest(
            data.dropna(),
            'norm',
            args=(data.mean(), data.std())
        )

        # Skewness and kurtosis
        skewness = stats.skew(data.dropna())
        kurt = stats.kurtosis(data.dropna())

        return {
            'shapiro_statistic': shapiro_stat,
            'shapiro_p_value': shapiro_p,
            'is_normal_shapiro': shapiro_p > 0.05 if shapiro_p else None,
            'ks_statistic': ks_stat,
            'ks_p_value': ks_p,
            'is_normal_ks': ks_p > 0.05,
            'skewness': skewness,
            'kurtosis': kurt,
            'anderson_statistic': anderson_result.statistic,
            'anderson_critical_values': anderson_result.critical_values
        }

    @staticmethod
    def compare_transformations(data: pd.DataFrame,
                               metric_col: str,
                               transformations: List[str] = None) -> pd.DataFrame:
        """
        Compare normality under different transformations.

        Parameters:
        -----------
        data : DataFrame
            Data to transform
        metric_col : str
            Column to transform
        transformations : List[str]
            List of transformations to try

        Returns:
        --------
        DataFrame comparing transformation results
        """
        if transformations is None:
            transformations = ['none', 'log', 'sqrt', 'boxcox']

        transformer = DataTransformer()
        results = []

        for trans in transformations:
            if trans == 'none':
                transformed = data[metric_col]
            elif trans == 'log':
                transformed = transformer.log_transform(data[metric_col])
            elif trans == 'sqrt':
                transformed = transformer.sqrt_transform(data[metric_col])
            elif trans == 'boxcox':
                transformed, _ = transformer.boxcox_transform(data[metric_col])
            else:
                continue

            # Assess normality
            normality = transformer.assess_normality(transformed)

            results.append({
                'transformation': trans,
                'skewness': normality['skewness'],
                'kurtosis': normality['kurtosis'],
                'shapiro_p': normality['shapiro_p_value'],
                'ks_p': normality['ks_p_value'],
                'is_normal': normality['is_normal_ks']
            })

        return pd.DataFrame(results).sort_values(
            'shapiro_p' if results[0]['shapiro_p'] else 'ks_p',
            ascending=False
        )
\end{lstlisting}

\subsection{Quality Assurance Checks}

Final QA before analysis:

\lstset{language=Python, caption=Pre-Analysis Quality Assurance}
\begin{lstlisting}
def run_pre_analysis_qa(data: pd.DataFrame,
                        variant_col: str = 'variant',
                        metric_cols: List[str] = None,
                        user_col: str = 'user_id',
                        timestamp_col: str = 'timestamp') -> Dict:
    """
    Comprehensive pre-analysis quality assurance.

    Returns:
    --------
    Dict with QA results and recommendations
    """
    qa_results = {
        'passed': True,
        'errors': [],
        'warnings': [],
        'recommendations': []
    }

    # 1. Data validation
    validator = ABTestDataValidator(
        data=data,
        variant_col=variant_col,
        user_col=user_col,
        timestamp_col=timestamp_col,
        metric_cols=metric_cols
    )

    validation = validator.run_all_checks()

    if not validation['summary']['passed']:
        qa_results['passed'] = False
        qa_results['errors'].extend(validation['summary']['issues'])

    qa_results['warnings'].extend(validation['summary']['warnings'])

    # 2. Missing data analysis
    if metric_cols:
        missing_handler = MissingDataHandler()
        missing_analysis = missing_handler.analyze_missingness(
            data[metric_cols]
        )

        high_missing = missing_analysis[
            missing_analysis['pct_missing'] > 5
        ]

        if len(high_missing) > 0:
            qa_results['warnings'].append(
                f"High missing data in {len(high_missing)} metrics"
            )
            qa_results['recommendations'].append(
                "Consider multiple imputation for metrics with >5% missing"
            )

    # 3. Outlier detection
    if metric_cols:
        detector = OutlierDetector()
        for metric in metric_cols:
            outliers_iqr = detector.detect_iqr(data[metric])
            pct_outliers = outliers_iqr.sum() / len(data) * 100

            if pct_outliers > 5:
                qa_results['warnings'].append(
                    f"{metric}: {pct_outliers:.1f}% outliers detected"
                )
                qa_results['recommendations'].append(
                    f"Consider winsorization for {metric}"
                )

    # 4. Normality assessment
    if metric_cols:
        transformer = DataTransformer()
        for metric in metric_cols:
            normality = transformer.assess_normality(data[metric].dropna())

            if not normality['is_normal_ks']:
                skew = normality['skewness']
                qa_results['recommendations'].append(
                    f"{metric} is non-normal (skewness={skew:.2f}). "
                    f"Consider transformation or non-parametric tests."
                )

    # 5. Sample size check
    variant_sizes = data[variant_col].value_counts()
    min_size = variant_sizes.min()

    if min_size < 100:
        qa_results['warnings'].append(
            f"Small sample size: minimum n={min_size}"
        )
        qa_results['recommendations'].append(
            "Small sample size may limit statistical power"
        )

    return qa_results

# Generate QA report
def generate_qa_report(qa_results: Dict) -> str:
    """Generate human-readable QA report."""
    report = []
    report.append("=" * 60)
    report.append("PRE-ANALYSIS QUALITY ASSURANCE REPORT")
    report.append("=" * 60)
    report.append("")

    # Status
    status = "PASSED" if qa_results['passed'] else "FAILED"
    report.append(f"Status: {status}")
    report.append("")

    # Errors
    if qa_results['errors']:
        report.append("ERRORS (Must Fix):")
        for error in qa_results['errors']:
            report.append(f"  X {error}")
        report.append("")

    # Warnings
    if qa_results['warnings']:
        report.append("WARNINGS:")
        for warning in qa_results['warnings']:
            report.append(f"  ! {warning}")
        report.append("")

    # Recommendations
    if qa_results['recommendations']:
        report.append("RECOMMENDATIONS:")
        for rec in qa_results['recommendations']:
            report.append(f"  - {rec}")
        report.append("")

    return "\n".join(report)
\end{lstlisting}

\section{Exploratory Data Analysis}

After preprocessing, conduct systematic exploratory analysis to understand data structure, distributions, and relationships.

\subsection{Descriptive Statistics}

Compute comprehensive summary statistics:

\lstset{language=Python, caption=Descriptive Statistics Framework}
\begin{lstlisting}
class DescriptiveStats:
    """Comprehensive descriptive statistics for A/B tests."""

    @staticmethod
    def summary_statistics(data: pd.DataFrame,
                          metric_col: str,
                          variant_col: str = 'variant') -> pd.DataFrame:
        """
        Calculate comprehensive summary statistics by variant.

        Returns:
        --------
        DataFrame with summary statistics
        """
        def compute_stats(series: pd.Series) -> Dict:
            """Compute statistics for a series."""
            return {
                'n': len(series),
                'n_missing': series.isnull().sum(),
                'mean': series.mean(),
                'median': series.median(),
                'std': series.std(),
                'min': series.min(),
                'q25': series.quantile(0.25),
                'q50': series.quantile(0.50),
                'q75': series.quantile(0.75),
                'max': series.max(),
                'iqr': series.quantile(0.75) - series.quantile(0.25),
                'skewness': stats.skew(series.dropna()),
                'kurtosis': stats.kurtosis(series.dropna()),
                'cv': series.std() / series.mean() if series.mean() != 0 else np.inf
            }

        # Overall statistics
        overall = compute_stats(data[metric_col])
        overall['variant'] = 'overall'

        # By variant
        variant_stats = []
        for variant in data[variant_col].unique():
            variant_data = data[data[variant_col] == variant][metric_col]
            stats_dict = compute_stats(variant_data)
            stats_dict['variant'] = variant
            variant_stats.append(stats_dict)

        # Combine
        all_stats = [overall] + variant_stats

        return pd.DataFrame(all_stats).set_index('variant')

    @staticmethod
    def confidence_intervals(data: pd.DataFrame,
                           metric_col: str,
                           variant_col: str = 'variant',
                           confidence_level: float = 0.95) -> pd.DataFrame:
        """
        Calculate confidence intervals for means.

        Returns:
        --------
        DataFrame with confidence intervals
        """
        results = []

        alpha = 1 - confidence_level

        for variant in data[variant_col].unique():
            variant_data = data[data[variant_col] == variant][metric_col].dropna()

            n = len(variant_data)
            mean = variant_data.mean()
            sem = stats.sem(variant_data)

            # T-distribution for CI
            ci = stats.t.interval(
                confidence_level,
                df=n-1,
                loc=mean,
                scale=sem
            )

            results.append({
                'variant': variant,
                'mean': mean,
                'sem': sem,
                'ci_lower': ci[0],
                'ci_upper': ci[1],
                'ci_width': ci[1] - ci[0],
                'n': n
            })

        return pd.DataFrame(results)

    @staticmethod
    def proportion_statistics(data: pd.DataFrame,
                             success_col: str,
                             variant_col: str = 'variant',
                             confidence_level: float = 0.95) -> pd.DataFrame:
        """
        Calculate statistics for binary/proportion metrics.

        Parameters:
        -----------
        data : DataFrame
            Experimental data
        success_col : str
            Binary column (0/1) indicating success
        variant_col : str
            Variant column
        confidence_level : float
            CI confidence level

        Returns:
        --------
        DataFrame with proportion statistics and CIs
        """
        from statsmodels.stats.proportion import proportion_confint

        results = []

        for variant in data[variant_col].unique():
            variant_data = data[data[variant_col] == variant]

            n = len(variant_data)
            successes = variant_data[success_col].sum()
            rate = successes / n

            # Wilson score interval (better for proportions)
            ci_lower, ci_upper = proportion_confint(
                successes,
                n,
                alpha=1-confidence_level,
                method='wilson'
            )

            # Standard error
            se = np.sqrt(rate * (1 - rate) / n)

            results.append({
                'variant': variant,
                'n': n,
                'successes': successes,
                'rate': rate,
                'se': se,
                'ci_lower': ci_lower,
                'ci_upper': ci_upper,
                'ci_width': ci_upper - ci_lower
            })

        return pd.DataFrame(results)
\end{lstlisting}

\subsection{Visualization Techniques}

Effective visualizations reveal patterns and communicate findings.

\lstset{language=Python, caption=Comprehensive EDA Visualizations}
\begin{lstlisting}
import matplotlib.pyplot as plt
import seaborn as sns
import plotly.graph_objects as go
import plotly.express as px
from plotly.subplots import make_subplots

class EDAVisualizer:
    """Create comprehensive EDA visualizations."""

    def __init__(self, style: str = 'seaborn-v0_8'):
        """Initialize visualizer with style."""
        plt.style.use(style)
        sns.set_palette("husl")

    def create_dashboard(self,
                        data: pd.DataFrame,
                        metric_col: str,
                        variant_col: str = 'variant',
                        timestamp_col: str = None) -> plt.Figure:
        """
        Create comprehensive EDA dashboard.

        Returns:
        --------
        Matplotlib figure with multiple subplots
        """
        fig = plt.figure(figsize=(16, 12))
        gs = fig.add_gridspec(3, 3, hspace=0.3, wspace=0.3)

        # 1. Distribution comparison (top left)
        ax1 = fig.add_subplot(gs[0, 0])
        for variant in data[variant_col].unique():
            variant_data = data[data[variant_col] == variant][metric_col].dropna()
            ax1.hist(variant_data, bins=50, alpha=0.6, label=f'{variant}')
        ax1.set_xlabel(metric_col)
        ax1.set_ylabel('Frequency')
        ax1.set_title('Distribution by Variant')
        ax1.legend()
        ax1.grid(alpha=0.3)

        # 2. Box plot comparison (top middle)
        ax2 = fig.add_subplot(gs[0, 1])
        data.boxplot(column=metric_col, by=variant_col, ax=ax2)
        ax2.set_xlabel('Variant')
        ax2.set_ylabel(metric_col)
        ax2.set_title('Box Plot Comparison')
        plt.sca(ax2)
        plt.xticks(rotation=0)

        # 3. Violin plot (top right)
        ax3 = fig.add_subplot(gs[0, 2])
        sns.violinplot(data=data, x=variant_col, y=metric_col, ax=ax3)
        ax3.set_title('Violin Plot')
        ax3.grid(alpha=0.3)

        # 4. Mean with confidence intervals (middle left)
        ax4 = fig.add_subplot(gs[1, 0])
        desc_stats = DescriptiveStats()
        ci_data = desc_stats.confidence_intervals(data, metric_col, variant_col)

        x_pos = range(len(ci_data))
        ax4.errorbar(
            x_pos,
            ci_data['mean'],
            yerr=[
                ci_data['mean'] - ci_data['ci_lower'],
                ci_data['ci_upper'] - ci_data['mean']
            ],
            fmt='o',
            markersize=8,
            capsize=5,
            capthick=2
        )
        ax4.set_xticks(x_pos)
        ax4.set_xticklabels(ci_data['variant'])
        ax4.set_ylabel(f'Mean {metric_col}')
        ax4.set_title('Means with 95% Confidence Intervals')
        ax4.grid(alpha=0.3)

        # 5. Q-Q plots (middle middle and right)
        variants = data[variant_col].unique()
        for idx, variant in enumerate(variants[:2]):
            ax = fig.add_subplot(gs[1, idx + 1])
            variant_data = data[data[variant_col] == variant][metric_col].dropna()
            stats.probplot(variant_data, dist="norm", plot=ax)
            ax.set_title(f'Q-Q Plot: {variant}')
            ax.grid(alpha=0.3)

        # 6. Time series (bottom row - if timestamp available)
        if timestamp_col:
            ax6 = fig.add_subplot(gs[2, :])
            data[timestamp_col] = pd.to_datetime(data[timestamp_col])

            # Daily aggregation
            daily_data = data.groupby(
                [pd.Grouper(key=timestamp_col, freq='D'), variant_col]
            )[metric_col].agg(['mean', 'sem', 'count']).reset_index()

            for variant in data[variant_col].unique():
                variant_daily = daily_data[daily_data[variant_col] == variant]

                ax6.plot(
                    variant_daily[timestamp_col],
                    variant_daily['mean'],
                    marker='o',
                    label=f'{variant}',
                    linewidth=2
                )

                # Add confidence bands
                ax6.fill_between(
                    variant_daily[timestamp_col],
                    variant_daily['mean'] - 1.96 * variant_daily['sem'],
                    variant_daily['mean'] + 1.96 * variant_daily['sem'],
                    alpha=0.2
                )

            ax6.set_xlabel('Date')
            ax6.set_ylabel(f'Daily Mean {metric_col}')
            ax6.set_title('Time Series with 95% Confidence Bands')
            ax6.legend()
            ax6.grid(alpha=0.3)
            plt.setp(ax6.xaxis.get_majorticklabels(), rotation=45)
        else:
            # 7. CDF comparison (bottom left)
            ax7 = fig.add_subplot(gs[2, 0])
            for variant in data[variant_col].unique():
                variant_data = data[
                    data[variant_col] == variant
                ][metric_col].dropna().sort_values()

                cdf = np.arange(1, len(variant_data) + 1) / len(variant_data)
                ax7.plot(variant_data, cdf, label=f'{variant}', linewidth=2)

            ax7.set_xlabel(metric_col)
            ax7.set_ylabel('Cumulative Probability')
            ax7.set_title('Empirical CDF Comparison')
            ax7.legend()
            ax7.grid(alpha=0.3)

            # 8. Density plot (bottom middle)
            ax8 = fig.add_subplot(gs[2, 1])
            for variant in data[variant_col].unique():
                variant_data = data[data[variant_col] == variant][metric_col].dropna()
                sns.kdeplot(data=variant_data, label=f'{variant}', ax=ax8, linewidth=2)
            ax8.set_xlabel(metric_col)
            ax8.set_ylabel('Density')
            ax8.set_title('Kernel Density Estimation')
            ax8.legend()
            ax8.grid(alpha=0.3)

            # 9. Summary statistics table (bottom right)
            ax9 = fig.add_subplot(gs[2, 2])
            ax9.axis('tight')
            ax9.axis('off')

            summary = desc_stats.summary_statistics(data, metric_col, variant_col)
            summary_display = summary[['mean', 'median', 'std', 'n']].round(2)

            table = ax9.table(
                cellText=summary_display.values,
                colLabels=summary_display.columns,
                rowLabels=summary_display.index,
                cellLoc='center',
                loc='center'
            )
            table.auto_set_font_size(False)
            table.set_fontsize(9)
            table.scale(1, 2)
            ax9.set_title('Summary Statistics')

        plt.suptitle(f'Exploratory Data Analysis: {metric_col}',
                     fontsize=16, y=0.995)

        return fig

    def create_interactive_dashboard(self,
                                    data: pd.DataFrame,
                                    metric_col: str,
                                    variant_col: str = 'variant',
                                    timestamp_col: str = None):
        """
        Create interactive Plotly dashboard.

        Returns:
        --------
        Plotly figure object
        """
        # Create subplots
        fig = make_subplots(
            rows=2, cols=2,
            subplot_titles=(
                'Distribution Comparison',
                'Box Plot',
                'Time Series' if timestamp_col else 'CDF Comparison',
                'Mean Comparison'
            ),
            specs=[
                [{"type": "histogram"}, {"type": "box"}],
                [{"type": "scatter"}, {"type": "bar"}]
            ]
        )

        # 1. Histogram
        for variant in data[variant_col].unique():
            variant_data = data[data[variant_col] == variant][metric_col].dropna()
            fig.add_trace(
                go.Histogram(
                    x=variant_data,
                    name=str(variant),
                    opacity=0.7,
                    nbinsx=50
                ),
                row=1, col=1
            )

        # 2. Box plot
        for variant in data[variant_col].unique():
            variant_data = data[data[variant_col] == variant][metric_col].dropna()
            fig.add_trace(
                go.Box(
                    y=variant_data,
                    name=str(variant),
                    boxmean='sd'
                ),
                row=1, col=2
            )

        # 3. Time series or CDF
        if timestamp_col:
            data[timestamp_col] = pd.to_datetime(data[timestamp_col])
            daily_data = data.groupby(
                [pd.Grouper(key=timestamp_col, freq='D'), variant_col]
            )[metric_col].mean().reset_index()

            for variant in data[variant_col].unique():
                variant_daily = daily_data[daily_data[variant_col] == variant]
                fig.add_trace(
                    go.Scatter(
                        x=variant_daily[timestamp_col],
                        y=variant_daily[metric_col],
                        name=str(variant),
                        mode='lines+markers'
                    ),
                    row=2, col=1
                )
        else:
            # CDF
            for variant in data[variant_col].unique():
                variant_data = data[
                    data[variant_col] == variant
                ][metric_col].dropna().sort_values()

                cdf = np.arange(1, len(variant_data) + 1) / len(variant_data)
                fig.add_trace(
                    go.Scatter(
                        x=variant_data,
                        y=cdf,
                        name=str(variant),
                        mode='lines'
                    ),
                    row=2, col=1
                )

        # 4. Mean comparison
        desc_stats = DescriptiveStats()
        ci_data = desc_stats.confidence_intervals(data, metric_col, variant_col)

        fig.add_trace(
            go.Bar(
                x=ci_data['variant'].astype(str),
                y=ci_data['mean'],
                error_y=dict(
                    type='data',
                    symmetric=False,
                    array=ci_data['ci_upper'] - ci_data['mean'],
                    arrayminus=ci_data['mean'] - ci_data['ci_lower']
                ),
                name='Mean with 95% CI'
            ),
            row=2, col=2
        )

        # Update layout
        fig.update_layout(
            height=800,
            showlegend=True,
            title_text=f"Interactive EDA Dashboard: {metric_col}"
        )

        fig.update_xaxes(title_text=metric_col, row=1, col=1)
        fig.update_yaxes(title_text="Frequency", row=1, col=1)

        fig.update_yaxes(title_text=metric_col, row=1, col=2)

        if timestamp_col:
            fig.update_xaxes(title_text="Date", row=2, col=1)
            fig.update_yaxes(title_text=f"Mean {metric_col}", row=2, col=1)
        else:
            fig.update_xaxes(title_text=metric_col, row=2, col=1)
            fig.update_yaxes(title_text="Cumulative Probability", row=2, col=1)

        fig.update_xaxes(title_text="Variant", row=2, col=2)
        fig.update_yaxes(title_text=f"Mean {metric_col}", row=2, col=2)

        return fig
\end{lstlisting}

\subsection{Distribution Analysis}

Understanding the underlying distribution of metrics is critical for selecting appropriate statistical tests and interpreting results.

\lstset{language=Python, caption=Distribution Analysis Functions}
\begin{lstlisting}
class DistributionAnalyzer:
    """Analyze and compare distributions in A/B test data."""

    @staticmethod
    def compare_distributions(data: pd.DataFrame,
                            metric_col: str,
                            variant_col: str = 'variant') -> Dict:
        """
        Compare distributions between variants using statistical tests.

        Returns:
        --------
        Dict with distribution comparison results
        """
        from scipy import stats

        variants = data[variant_col].unique()
        if len(variants) != 2:
            raise ValueError("This function requires exactly 2 variants")

        control = variants[0]
        treatment = variants[1]

        control_data = data[data[variant_col] == control][metric_col].dropna()
        treatment_data = data[data[variant_col] == treatment][metric_col].dropna()

        results = {}

        # Kolmogorov-Smirnov test
        ks_stat, ks_p = stats.ks_2samp(control_data, treatment_data)
        results['ks_test'] = {
            'statistic': ks_stat,
            'p_value': ks_p,
            'distributions_differ': ks_p < 0.05
        }

        # Mann-Whitney U test (for medians)
        mw_stat, mw_p = stats.mannwhitneyu(
            control_data, treatment_data, alternative='two-sided'
        )
        results['mann_whitney'] = {
            'statistic': mw_stat,
            'p_value': mw_p,
            'medians_differ': mw_p < 0.05
        }

        # Levene's test (for variance equality)
        levene_stat, levene_p = stats.levene(control_data, treatment_data)
        results['levene_test'] = {
            'statistic': levene_stat,
            'p_value': levene_p,
            'variances_differ': levene_p < 0.05,
            'homoscedastic': levene_p >= 0.05
        }

        return results

    @staticmethod
    def plot_distribution_comparison(data: pd.DataFrame,
                                    metric_col: str,
                                    variant_col: str = 'variant'):
        """Create detailed distribution comparison visualization."""
        import matplotlib.pyplot as plt
        import seaborn as sns

        fig, axes = plt.subplots(2, 2, figsize=(14, 10))

        # 1. Overlapping histograms
        for variant in data[variant_col].unique():
            variant_data = data[data[variant_col] == variant][metric_col].dropna()
            axes[0, 0].hist(variant_data, bins=50, alpha=0.6, label=str(variant))
        axes[0, 0].set_xlabel(metric_col)
        axes[0, 0].set_ylabel('Frequency')
        axes[0, 0].set_title('Histogram Comparison')
        axes[0, 0].legend()
        axes[0, 0].grid(alpha=0.3)

        # 2. KDE plots
        for variant in data[variant_col].unique():
            variant_data = data[data[variant_col] == variant][metric_col].dropna()
            sns.kdeplot(data=variant_data, label=str(variant),
                       ax=axes[0, 1], linewidth=2)
        axes[0, 1].set_xlabel(metric_col)
        axes[0, 1].set_ylabel('Density')
        axes[0, 1].set_title('Kernel Density Estimation')
        axes[0, 1].legend()
        axes[0, 1].grid(alpha=0.3)

        # 3. Empirical CDF
        for variant in data[variant_col].unique():
            variant_data = data[
                data[variant_col] == variant
            ][metric_col].dropna().sort_values()
            cdf = np.arange(1, len(variant_data) + 1) / len(variant_data)
            axes[1, 0].plot(variant_data, cdf, label=str(variant), linewidth=2)
        axes[1, 0].set_xlabel(metric_col)
        axes[1, 0].set_ylabel('Cumulative Probability')
        axes[1, 0].set_title('Empirical CDF')
        axes[1, 0].legend()
        axes[1, 0].grid(alpha=0.3)

        # 4. Box and violin plots combined
        ax_box = axes[1, 1]
        positions = np.arange(len(data[variant_col].unique()))

        # Violin plots
        parts = ax_box.violinplot(
            [data[data[variant_col] == v][metric_col].dropna()
             for v in data[variant_col].unique()],
            positions=positions,
            showmeans=False,
            showmedians=False
        )

        # Overlay box plots
        ax_box.boxplot(
            [data[data[variant_col] == v][metric_col].dropna()
             for v in data[variant_col].unique()],
            positions=positions,
            widths=0.15
        )

        ax_box.set_xticks(positions)
        ax_box.set_xticklabels(data[variant_col].unique())
        ax_box.set_ylabel(metric_col)
        ax_box.set_title('Box + Violin Plot')
        ax_box.grid(alpha=0.3)

        plt.tight_layout()
        return fig

\end{lstlisting}

\subsection{Trend and Seasonality Detection}

Temporal patterns can indicate novelty effects, external events, or seasonal variations.

\lstset{language=Python, caption=Temporal Pattern Detection}
\begin{lstlisting}
class TemporalAnalyzer:
    """Analyze temporal patterns in A/B test data."""

    @staticmethod
    def detect_trend(data: pd.DataFrame,
                    metric_col: str,
                    timestamp_col: str,
                    variant_col: str = None) -> Dict:
        """
        Detect linear trend in time series using Mann-Kendall test.

        Returns:
        --------
        Dict with trend detection results
        """
        import pymannkendall as mk

        data = data.copy()
        data[timestamp_col] = pd.to_datetime(data[timestamp_col])
        data = data.sort_values(timestamp_col)

        if variant_col:
            # Test within each variant
            results = {}
            for variant in data[variant_col].unique():
                variant_data = data[data[variant_col] == variant]
                daily_mean = variant_data.groupby(
                    variant_data[timestamp_col].dt.date
                )[metric_col].mean()

                mk_result = mk.original_test(daily_mean)

                results[variant] = {
                    'trend': mk_result.trend,
                    'p_value': mk_result.p,
                    'tau': mk_result.Tau,
                    'has_trend': mk_result.trend != 'no trend'
                }
            return results
        else:
            # Overall trend
            daily_mean = data.groupby(
                data[timestamp_col].dt.date
            )[metric_col].mean()

            mk_result = mk.original_test(daily_mean)

            return {
                'trend': mk_result.trend,
                'p_value': mk_result.p,
                'tau': mk_result.Tau,
                'has_trend': mk_result.trend != 'no trend'
            }

    @staticmethod
    def detect_day_of_week_effects(data: pd.DataFrame,
                                   metric_col: str,
                                   timestamp_col: str,
                                   variant_col: str = 'variant') -> Dict:
        """
        Test for day-of-week effects using ANOVA.

        Returns:
        --------
        Dict with day-of-week analysis results
        """
        from scipy import stats as sp_stats

        data = data.copy()
        data[timestamp_col] = pd.to_datetime(data[timestamp_col])
        data['day_of_week'] = data[timestamp_col].dt.dayofweek
        data['day_name'] = data[timestamp_col].dt.day_name()

        results = {}

        # Overall effect
        groups = [data[data['day_of_week'] == d][metric_col].dropna()
                 for d in range(7)]
        f_stat, p_value = sp_stats.f_oneway(*groups)

        results['overall'] = {
            'f_statistic': f_stat,
            'p_value': p_value,
            'has_dow_effect': p_value < 0.05
        }

        # By variant
        variant_results = {}
        for variant in data[variant_col].unique():
            variant_data = data[data[variant_col] == variant]
            groups = [variant_data[variant_data['day_of_week'] == d][metric_col].dropna()
                     for d in range(7)]
            f_stat, p_value = sp_stats.f_oneway(*groups)

            variant_results[variant] = {
                'f_statistic': f_stat,
                'p_value': p_value,
                'has_dow_effect': p_value < 0.05
            }

        results['by_variant'] = variant_results

        # Mean by day of week
        dow_means = data.groupby(['day_name', variant_col])[metric_col].mean().unstack()
        results['means_by_dow'] = dow_means

        return results

    @staticmethod
    def plot_temporal_patterns(data: pd.DataFrame,
                              metric_col: str,
                              timestamp_col: str,
                              variant_col: str = 'variant'):
        """Create comprehensive temporal analysis visualizations."""
        import matplotlib.pyplot as plt

        data = data.copy()
        data[timestamp_col] = pd.to_datetime(data[timestamp_col])
        data['date'] = data[timestamp_col].dt.date
        data['hour'] = data[timestamp_col].dt.hour
        data['day_of_week'] = data[timestamp_col].dt.dayofweek
        data['day_name'] = data[timestamp_col].dt.day_name()

        fig, axes = plt.subplots(2, 2, figsize=(16, 10))

        # 1. Daily trend
        daily_data = data.groupby(['date', variant_col])[metric_col].mean().unstack()

        for variant in daily_data.columns:
            axes[0, 0].plot(daily_data.index, daily_data[variant],
                          marker='o', label=str(variant), linewidth=2, alpha=0.7)

        axes[0, 0].set_xlabel('Date')
        axes[0, 0].set_ylabel(f'Mean {metric_col}')
        axes[0, 0].set_title('Daily Trend')
        axes[0, 0].legend()
        axes[0, 0].grid(alpha=0.3)
        axes[0, 0].tick_params(axis='x', rotation=45)

        # 2. Day of week pattern
        dow_order = ['Monday', 'Tuesday', 'Wednesday', 'Thursday',
                     'Friday', 'Saturday', 'Sunday']
        dow_data = data.groupby(['day_name', variant_col])[metric_col].mean().unstack()
        dow_data = dow_data.reindex(dow_order)

        x = np.arange(len(dow_order))
        width = 0.35

        variants = data[variant_col].unique()
        for i, variant in enumerate(variants):
            offset = width * (i - len(variants)/2 + 0.5)
            axes[0, 1].bar(x + offset, dow_data[variant],
                          width, label=str(variant), alpha=0.7)

        axes[0, 1].set_xlabel('Day of Week')
        axes[0, 1].set_ylabel(f'Mean {metric_col}')
        axes[0, 1].set_title('Day of Week Pattern')
        axes[0, 1].set_xticks(x)
        axes[0, 1].set_xticklabels(dow_order, rotation=45)
        axes[0, 1].legend()
        axes[0, 1].grid(alpha=0.3)

        # 3. Hourly pattern (if timestamps include hours)
        if data['hour'].nunique() > 1:
            hourly_data = data.groupby(['hour', variant_col])[metric_col].mean().unstack()

            for variant in hourly_data.columns:
                axes[1, 0].plot(hourly_data.index, hourly_data[variant],
                              marker='o', label=str(variant), linewidth=2)

            axes[1, 0].set_xlabel('Hour of Day')
            axes[1, 0].set_ylabel(f'Mean {metric_col}')
            axes[1, 0].set_title('Hourly Pattern')
            axes[1, 0].legend()
            axes[1, 0].grid(alpha=0.3)

        # 4. Cumulative metric over time
        cumulative_data = data.sort_values(timestamp_col).groupby(
            variant_col
        )[metric_col].cumsum().reset_index()

        for variant in data[variant_col].unique():
            variant_data = data[data[variant_col] == variant].sort_values(timestamp_col)
            variant_cumulative = variant_data[metric_col].cumsum()

            axes[1, 1].plot(range(len(variant_cumulative)), variant_cumulative,
                          label=str(variant), linewidth=2, alpha=0.7)

        axes[1, 1].set_xlabel('Observation Index')
        axes[1, 1].set_ylabel(f'Cumulative {metric_col}')
        axes[1, 1].set_title('Cumulative Metric Over Time')
        axes[1, 1].legend()
        axes[1, 1].grid(alpha=0.3)

        plt.tight_layout()
        return fig

\end{lstlisting}

\section{Statistical Analysis Pipeline}

After preprocessing and EDA, execute the statistical tests systematically.

\subsection{Assumption Checking}

Before applying parametric tests, verify assumptions.

\lstset{language=Python, caption=Statistical Assumption Checking}
\begin{lstlisting}
class AssumptionChecker:
    """Check statistical assumptions for A/B test analysis."""

    @staticmethod
    def check_normality(data: pd.DataFrame,
                       metric_col: str,
                       variant_col: str = 'variant',
                       alpha: float = 0.05) -> Dict:
        """
        Check normality assumption for each variant.

        Returns:
        --------
        Dict with normality test results
        """
        results = {}

        for variant in data[variant_col].unique():
            variant_data = data[data[variant_col] == variant][metric_col].dropna()

            # Shapiro-Wilk test
            if len(variant_data) < 5000:
                sw_stat, sw_p = stats.shapiro(variant_data)
            else:
                sw_stat, sw_p = None, None

            # Anderson-Darling test
            ad_result = stats.anderson(variant_data)

            # K-S test
            ks_stat, ks_p = stats.kstest(
                variant_data,
                'norm',
                args=(variant_data.mean(), variant_data.std())
            )

            results[variant] = {
                'shapiro_wilk': {
                    'statistic': sw_stat,
                    'p_value': sw_p,
                    'is_normal': sw_p > alpha if sw_p else None
                },
                'kolmogorov_smirnov': {
                    'statistic': ks_stat,
                    'p_value': ks_p,
                    'is_normal': ks_p > alpha
                },
                'anderson_darling': {
                    'statistic': ad_result.statistic,
                    'critical_values': ad_result.critical_values.tolist(),
                    'significance_levels': ad_result.significance_level.tolist()
                },
                'skewness': stats.skew(variant_data),
                'kurtosis': stats.kurtosis(variant_data)
            }

        return results

    @staticmethod
    def check_variance_equality(data: pd.DataFrame,
                               metric_col: str,
                               variant_col: str = 'variant',
                               alpha: float = 0.05) -> Dict:
        """
        Check variance equality assumption (homoscedasticity).

        Returns:
        --------
        Dict with variance equality test results
        """
        variants = data[variant_col].unique()
        groups = [data[data[variant_col] == v][metric_col].dropna()
                 for v in variants]

        # Levene's test (robust to non-normality)
        levene_stat, levene_p = stats.levene(*groups)

        # Bartlett's test (assumes normality)
        bartlett_stat, bartlett_p = stats.bartlett(*groups)

        # Fligner-Killeen test (non-parametric)
        fligner_stat, fligner_p = stats.fligner(*groups)

        # Variance ratios
        variances = {v: data[data[variant_col] == v][metric_col].var()
                    for v in variants}

        return {
            'levene_test': {
                'statistic': levene_stat,
                'p_value': levene_p,
                'equal_variances': levene_p > alpha
            },
            'bartlett_test': {
                'statistic': bartlett_stat,
                'p_value': bartlett_p,
                'equal_variances': bartlett_p > alpha
            },
            'fligner_test': {
                'statistic': fligner_stat,
                'p_value': fligner_p,
                'equal_variances': fligner_p > alpha
            },
            'variances_by_variant': variances,
            'variance_ratio': max(variances.values()) / min(variances.values())
        }

    @staticmethod
    def check_independence(data: pd.DataFrame,
                          user_col: str = 'user_id',
                          variant_col: str = 'variant') -> Dict:
        """
        Check independence assumption.

        Returns:
        --------
        Dict with independence check results
        """
        # Check for users appearing in multiple variants
        users_per_variant = data.groupby(user_col)[variant_col].nunique()
        multi_variant_users = users_per_variant[users_per_variant > 1]

        # Check for temporal autocorrelation (if timestamp available)
        results = {
            'multi_variant_users': {
                'count': len(multi_variant_users),
                'percentage': len(multi_variant_users) / len(users_per_variant) * 100,
                'independent': len(multi_variant_users) == 0
            }
        }

        return results

    @staticmethod
    def generate_assumption_report(data: pd.DataFrame,
                                   metric_col: str,
                                   variant_col: str = 'variant',
                                   user_col: str = 'user_id') -> str:
        """Generate comprehensive assumption check report."""
        checker = AssumptionChecker()

        report = []
        report.append("=" * 60)
        report.append("STATISTICAL ASSUMPTIONS CHECK")
        report.append("=" * 60)
        report.append("")

        # Normality
        report.append("1. NORMALITY:")
        normality = checker.check_normality(data, metric_col, variant_col)
        for variant, results in normality.items():
            report.append(f"\n  Variant: {variant}")
            if results['shapiro_wilk']['p_value']:
                report.append(f"    Shapiro-Wilk p-value: {results['shapiro_wilk']['p_value']:.4f}")
            report.append(f"    K-S p-value: {results['kolmogorov_smirnov']['p_value']:.4f}")
            report.append(f"    Skewness: {results['skewness']:.4f}")
            report.append(f"    Kurtosis: {results['kurtosis']:.4f}")

        # Variance equality
        report.append("\n2. VARIANCE EQUALITY:")
        var_eq = checker.check_variance_equality(data, metric_col, variant_col)
        report.append(f"  Levene's test p-value: {var_eq['levene_test']['p_value']:.4f}")
        report.append(f"  Equal variances: {var_eq['levene_test']['equal_variances']}")
        report.append(f"  Variance ratio: {var_eq['variance_ratio']:.2f}")

        # Independence
        report.append("\n3. INDEPENDENCE:")
        indep = checker.check_independence(data, user_col, variant_col)
        report.append(f"  Multi-variant users: {indep['multi_variant_users']['count']}")
        report.append(f"  Independent: {indep['multi_variant_users']['independent']}")

        report.append("")
        return "\n".join(report)

\end{lstlisting}

\subsection{Test Execution}

Execute appropriate statistical tests based on data characteristics.

\lstset{language=Python, caption=Automated Test Execution}
\begin{lstlisting}
class ABTestExecutor:
    """Execute statistical tests for A/B experiments."""

    @staticmethod
    def t_test(data: pd.DataFrame,
              metric_col: str,
              variant_col: str = 'variant',
              equal_var: bool = None,
              alpha: float = 0.05) -> Dict:
        """
        Execute t-test for continuous metrics.

        Parameters:
        -----------
        data : DataFrame
            Experimental data
        metric_col : str
            Metric to test
        variant_col : str
            Variant column
        equal_var : bool
            Assume equal variance? If None, will test
        alpha : float
            Significance level

        Returns:
        --------
        Dict with t-test results
        """
        variants = data[variant_col].unique()
        if len(variants) != 2:
            raise ValueError("T-test requires exactly 2 variants")

        control = variants[0]
        treatment = variants[1]

        control_data = data[data[variant_col] == control][metric_col].dropna()
        treatment_data = data[data[variant_col] == treatment][metric_col].dropna()

        # Test for equal variance if not specified
        if equal_var is None:
            _, levene_p = stats.levene(control_data, treatment_data)
            equal_var = levene_p > 0.05

        # Perform t-test
        t_stat, p_value = stats.ttest_ind(
            treatment_data,
            control_data,
            equal_var=equal_var
        )

        # Calculate effect size (Cohen's d)
        pooled_std = np.sqrt(
            ((len(control_data) - 1) * control_data.var() +
             (len(treatment_data) - 1) * treatment_data.var()) /
            (len(control_data) + len(treatment_data) - 2)
        )
        cohens_d = (treatment_data.mean() - control_data.mean()) / pooled_std

        # Confidence interval for difference
        se_diff = pooled_std * np.sqrt(
            1/len(control_data) + 1/len(treatment_data)
        )
        df = len(control_data) + len(treatment_data) - 2
        t_crit = stats.t.ppf(1 - alpha/2, df)

        diff = treatment_data.mean() - control_data.mean()
        ci_lower = diff - t_crit * se_diff
        ci_upper = diff + t_crit * se_diff

        return {
            'test_type': 't-test',
            'control_mean': control_data.mean(),
            'treatment_mean': treatment_data.mean(),
            'difference': diff,
            'relative_difference_pct': (diff / control_data.mean() * 100)
                                      if control_data.mean() != 0 else np.inf,
            't_statistic': t_stat,
            'p_value': p_value,
            'significant': p_value < alpha,
            'degrees_of_freedom': df,
            'equal_variance_assumed': equal_var,
            'cohens_d': cohens_d,
            'confidence_interval': {
                'level': 1 - alpha,
                'lower': ci_lower,
                'upper': ci_upper
            },
            'sample_sizes': {
                'control': len(control_data),
                'treatment': len(treatment_data)
            }
        }

    @staticmethod
    def mann_whitney_test(data: pd.DataFrame,
                         metric_col: str,
                         variant_col: str = 'variant',
                         alpha: float = 0.05) -> Dict:
        """
        Execute Mann-Whitney U test (non-parametric).

        Returns:
        --------
        Dict with Mann-Whitney test results
        """
        variants = data[variant_col].unique()
        if len(variants) != 2:
            raise ValueError("Mann-Whitney test requires exactly 2 variants")

        control = variants[0]
        treatment = variants[1]

        control_data = data[data[variant_col] == control][metric_col].dropna()
        treatment_data = data[data[variant_col] == treatment][metric_col].dropna()

        # Perform test
        u_stat, p_value = stats.mannwhitneyu(
            treatment_data,
            control_data,
            alternative='two-sided'
        )

        # Effect size (rank-biserial correlation)
        n1 = len(control_data)
        n2 = len(treatment_data)
        rank_biserial = 1 - (2*u_stat) / (n1 * n2)

        return {
            'test_type': 'mann_whitney',
            'control_median': control_data.median(),
            'treatment_median': treatment_data.median(),
            'difference_in_medians': treatment_data.median() - control_data.median(),
            'u_statistic': u_stat,
            'p_value': p_value,
            'significant': p_value < alpha,
            'rank_biserial_correlation': rank_biserial,
            'sample_sizes': {
                'control': n1,
                'treatment': n2
            }
        }

    @staticmethod
    def proportion_test(data: pd.DataFrame,
                       success_col: str,
                       variant_col: str = 'variant',
                       alpha: float = 0.05) -> Dict:
        """
        Execute proportion test (z-test for proportions).

        Returns:
        --------
        Dict with proportion test results
        """
        from statsmodels.stats.proportion import proportions_ztest, proportion_confint

        variants = data[variant_col].unique()
        if len(variants) != 2:
            raise ValueError("Proportion test requires exactly 2 variants")

        control = variants[0]
        treatment = variants[1]

        control_data = data[data[variant_col] == control]
        treatment_data = data[data[variant_col] == treatment]

        n_control = len(control_data)
        n_treatment = len(treatment_data)

        successes_control = control_data[success_col].sum()
        successes_treatment = treatment_data[success_col].sum()

        p_control = successes_control / n_control
        p_treatment = successes_treatment / n_treatment

        # Z-test
        counts = np.array([successes_treatment, successes_control])
        nobs = np.array([n_treatment, n_control])

        z_stat, p_value = proportions_ztest(counts, nobs)

        # Confidence intervals
        ci_control = proportion_confint(
            successes_control, n_control, alpha=alpha, method='wilson'
        )
        ci_treatment = proportion_confint(
            successes_treatment, n_treatment, alpha=alpha, method='wilson'
        )

        # Relative difference (lift)
        relative_diff = (p_treatment - p_control) / p_control * 100 if p_control > 0 else np.inf

        # Absolute difference CI
        se_diff = np.sqrt(
            p_control * (1 - p_control) / n_control +
            p_treatment * (1 - p_treatment) / n_treatment
        )
        z_crit = stats.norm.ppf(1 - alpha/2)
        diff = p_treatment - p_control
        ci_diff_lower = diff - z_crit * se_diff
        ci_diff_upper = diff + z_crit * se_diff

        return {
            'test_type': 'proportion_ztest',
            'control_rate': p_control,
            'treatment_rate': p_treatment,
            'absolute_difference': diff,
            'relative_difference_pct': relative_diff,
            'z_statistic': z_stat,
            'p_value': p_value,
            'significant': p_value < alpha,
            'confidence_intervals': {
                'control': {'lower': ci_control[0], 'upper': ci_control[1]},
                'treatment': {'lower': ci_treatment[0], 'upper': ci_treatment[1]},
                'difference': {'lower': ci_diff_lower, 'upper': ci_diff_upper}
            },
            'sample_sizes': {
                'control': n_control,
                'treatment': n_treatment
            },
            'successes': {
                'control': int(successes_control),
                'treatment': int(successes_treatment)
            }
        }

    @staticmethod
    def automatic_test_selection(data: pd.DataFrame,
                                 metric_col: str,
                                 variant_col: str = 'variant',
                                 is_binary: bool = None,
                                 alpha: float = 0.05) -> Dict:
        """
        Automatically select and execute appropriate test.

        Returns:
        --------
        Dict with test results and recommendation
        """
        executor = ABTestExecutor()
        checker = AssumptionChecker()

        # Determine if binary
        if is_binary is None:
            unique_vals = data[metric_col].dropna().unique()
            is_binary = len(unique_vals) == 2 and set(unique_vals) <= {0, 1}

        if is_binary:
            # Use proportion test
            result = executor.proportion_test(data, metric_col, variant_col, alpha)
            result['recommendation'] = "Proportion z-test (binary outcome)"
            return result

        # Check normality and variance equality
        normality = checker.check_normality(data, metric_col, variant_col, alpha)
        var_equality = checker.check_variance_equality(data, metric_col, variant_col, alpha)

        # Decision logic
        all_normal = all(
            v.get('shapiro_wilk', {}).get('is_normal', False) or
            v.get('kolmogorov_smirnov', {}).get('is_normal', False)
            for v in normality.values()
        )

        equal_var = var_equality['levene_test']['equal_variances']

        if all_normal:
            # Use t-test
            result = executor.t_test(data, metric_col, variant_col, equal_var, alpha)
            if equal_var:
                result['recommendation'] = "Student's t-test (equal variances)"
            else:
                result['recommendation'] = "Welch's t-test (unequal variances)"
        else:
            # Use Mann-Whitney
            result = executor.mann_whitney_test(data, metric_col, variant_col, alpha)
            result['recommendation'] = "Mann-Whitney U test (non-normal distribution)"

        result['assumption_checks'] = {
            'normality': normality,
            'variance_equality': var_equality
        }

        return result

\end{lstlisting}

\subsection{Effect Size Calculation}

Statistical significance doesn't indicate practical importance. Always calculate and report effect sizes.

\lstset{language=Python, caption=Effect Size Calculations}
\begin{lstlisting}
class EffectSizeCalculator:
    """Calculate various effect size measures for A/B tests."""

    @staticmethod
    def cohens_d(data: pd.DataFrame,
                metric_col: str,
                variant_col: str = 'variant',
                pooled: bool = True) -> Dict:
        """
        Calculate Cohen's d for continuous metrics.

        Parameters:
        -----------
        data : DataFrame
            Experimental data
        metric_col : str
            Metric to analyze
        variant_col : str
            Variant column
        pooled : bool
            Use pooled standard deviation

        Returns:
        --------
        Dict with Cohen's d and interpretation
        """
        variants = data[variant_col].unique()
        if len(variants) != 2:
            raise ValueError("Cohen's d requires exactly 2 variants")

        control = variants[0]
        treatment = variants[1]

        control_data = data[data[variant_col] == control][metric_col].dropna()
        treatment_data = data[data[variant_col] == treatment][metric_col].dropna()

        mean_diff = treatment_data.mean() - control_data.mean()

        if pooled:
            # Pooled standard deviation
            n1, n2 = len(control_data), len(treatment_data)
            var1, var2 = control_data.var(), treatment_data.var()
            pooled_std = np.sqrt(((n1 - 1) * var1 + (n2 - 1) * var2) / (n1 + n2 - 2))
            d = mean_diff / pooled_std
        else:
            # Control group standard deviation
            d = mean_diff / control_data.std()

        # Interpretation
        if abs(d) < 0.2:
            interpretation = "negligible"
        elif abs(d) < 0.5:
            interpretation = "small"
        elif abs(d) < 0.8:
            interpretation = "medium"
        else:
            interpretation = "large"

        return {
            'cohens_d': d,
            'interpretation': interpretation,
            'mean_difference': mean_diff,
            'pooled_std': pooled_std if pooled else None
        }

    @staticmethod
    def relative_effect(data: pd.DataFrame,
                       metric_col: str,
                       variant_col: str = 'variant') -> Dict:
        """
        Calculate relative (percentage) effect.

        Returns:
        --------
        Dict with relative effect measures
        """
        variants = data[variant_col].unique()
        if len(variants) != 2:
            raise ValueError("Requires exactly 2 variants")

        control = variants[0]
        treatment = variants[1]

        control_mean = data[data[variant_col] == control][metric_col].mean()
        treatment_mean = data[data[variant_col] == treatment][metric_col].mean()

        absolute_diff = treatment_mean - control_mean
        relative_diff_pct = (absolute_diff / control_mean * 100) if control_mean != 0 else np.inf

        return {
            'control_mean': control_mean,
            'treatment_mean': treatment_mean,
            'absolute_difference': absolute_diff,
            'relative_difference_pct': relative_diff_pct,
            'lift_pct': relative_diff_pct
        }

    @staticmethod
    def odds_ratio(data: pd.DataFrame,
                  success_col: str,
                  variant_col: str = 'variant') -> Dict:
        """
        Calculate odds ratio for binary outcomes.

        Returns:
        --------
        Dict with odds ratio and confidence interval
        """
        from statsmodels.stats.contingency_tables import Table2x2

        variants = data[variant_col].unique()
        if len(variants) != 2:
            raise ValueError("Requires exactly 2 variants")

        control = variants[0]
        treatment = variants[1]

        # Create contingency table
        control_data = data[data[variant_col] == control]
        treatment_data = data[data[variant_col] == treatment]

        # [treatment_success, treatment_failure]
        # [control_success, control_failure]
        table = np.array([
            [treatment_data[success_col].sum(),
             len(treatment_data) - treatment_data[success_col].sum()],
            [control_data[success_col].sum(),
             len(control_data) - control_data[success_col].sum()]
        ])

        table2x2 = Table2x2(table)
        or_value = table2x2.oddsratio
        ci = table2x2.oddsratio_confint()

        # Interpretation
        if or_value > 1:
            interpretation = f"Treatment has {or_value:.2f}x odds of success"
        elif or_value < 1:
            interpretation = f"Control has {1/or_value:.2f}x odds of success"
        else:
            interpretation = "Equal odds"

        return {
            'odds_ratio': or_value,
            'confidence_interval': {'lower': ci[0], 'upper': ci[1]},
            'log_odds_ratio': np.log(or_value),
            'interpretation': interpretation
        }

    @staticmethod
    def risk_ratio(data: pd.DataFrame,
                  success_col: str,
                  variant_col: str = 'variant') -> Dict:
        """
        Calculate risk ratio (relative risk) for binary outcomes.

        Returns:
        --------
        Dict with risk ratio
        """
        variants = data[variant_col].unique()
        if len(variants) != 2:
            raise ValueError("Requires exactly 2 variants")

        control = variants[0]
        treatment = variants[1]

        control_data = data[data[variant_col] == control]
        treatment_data = data[data[variant_col] == treatment]

        p_control = control_data[success_col].mean()
        p_treatment = treatment_data[success_col].mean()

        rr = p_treatment / p_control if p_control > 0 else np.inf

        # Confidence interval (log method)
        se_log_rr = np.sqrt(
            (1 - p_control) / (control_data[success_col].sum()) +
            (1 - p_treatment) / (treatment_data[success_col].sum())
        )

        ci_lower = np.exp(np.log(rr) - 1.96 * se_log_rr)
        ci_upper = np.exp(np.log(rr) + 1.96 * se_log_rr)

        return {
            'risk_ratio': rr,
            'confidence_interval': {'lower': ci_lower, 'upper': ci_upper},
            'control_risk': p_control,
            'treatment_risk': p_treatment
        }

\end{lstlisting}

\subsection{Confidence Interval Estimation}

Confidence intervals quantify uncertainty and are more informative than p-values alone.

\lstset{language=Python, caption=Advanced Confidence Interval Methods}
\begin{lstlisting}
class ConfidenceIntervalEstimator:
    """Advanced confidence interval estimation methods."""

    @staticmethod
    def bootstrap_ci(data: pd.DataFrame,
                    metric_col: str,
                    variant_col: str = 'variant',
                    statistic: str = 'mean',
                    n_bootstrap: int = 10000,
                    confidence_level: float = 0.95) -> Dict:
        """
        Calculate bootstrap confidence intervals.

        Parameters:
        -----------
        data : DataFrame
            Experimental data
        metric_col : str
            Metric to analyze
        variant_col : str
            Variant column
        statistic : str
            Statistic to compute: 'mean', 'median', 'std'
        n_bootstrap : int
            Number of bootstrap samples
        confidence_level : float
            Confidence level

        Returns:
        --------
        Dict with bootstrap CI
        """
        from scipy.stats import bootstrap

        variants = data[variant_col].unique()
        if len(variants) != 2:
            raise ValueError("Requires exactly 2 variants")

        control = variants[0]
        treatment = variants[1]

        control_data = data[data[variant_col] == control][metric_col].dropna().values
        treatment_data = data[data[variant_col] == treatment][metric_col].dropna().values

        # Define statistic function
        if statistic == 'mean':
            stat_func = np.mean
        elif statistic == 'median':
            stat_func = np.median
        elif statistic == 'std':
            stat_func = np.std
        else:
            raise ValueError(f"Unknown statistic: {statistic}")

        # Bootstrap for difference
        def difference_stat(x, y):
            return stat_func(y) - stat_func(x)

        bootstrap_diffs = []
        rng = np.random.RandomState(42)

        for _ in range(n_bootstrap):
            control_sample = rng.choice(control_data, size=len(control_data), replace=True)
            treatment_sample = rng.choice(treatment_data, size=len(treatment_data), replace=True)
            bootstrap_diffs.append(stat_func(treatment_sample) - stat_func(control_sample))

        bootstrap_diffs = np.array(bootstrap_diffs)

        # Percentile method
        alpha = 1 - confidence_level
        ci_lower = np.percentile(bootstrap_diffs, alpha/2 * 100)
        ci_upper = np.percentile(bootstrap_diffs, (1 - alpha/2) * 100)

        # Observed difference
        obs_diff = stat_func(treatment_data) - stat_func(control_data)

        return {
            'statistic': statistic,
            'observed_difference': obs_diff,
            'confidence_interval': {
                'method': 'percentile_bootstrap',
                'level': confidence_level,
                'lower': ci_lower,
                'upper': ci_upper
            },
            'n_bootstrap': n_bootstrap,
            'bootstrap_distribution': bootstrap_diffs
        }

    @staticmethod
    def delta_method_ci(data: pd.DataFrame,
                       metric_col: str,
                       variant_col: str = 'variant',
                       transformation: str = 'log',
                       confidence_level: float = 0.95) -> Dict:
        """
        Calculate CI using delta method for transformed statistics.

        Parameters:
        -----------
        transformation : str
            'log', 'logit', or 'sqrt'

        Returns:
        --------
        Dict with delta method CI
        """
        variants = data[variant_col].unique()
        if len(variants) != 2:
            raise ValueError("Requires exactly 2 variants")

        control = variants[0]
        treatment = variants[1]

        control_data = data[data[variant_col] == control][metric_col].dropna()
        treatment_data = data[data[variant_col] == treatment][metric_col].dropna()

        # Calculate means and variances
        mean_c = control_data.mean()
        mean_t = treatment_data.mean()
        var_c = control_data.var() / len(control_data)
        var_t = treatment_data.var() / len(treatment_data)

        alpha = 1 - confidence_level
        z_crit = stats.norm.ppf(1 - alpha/2)

        if transformation == 'log':
            # Log transformation
            ratio = mean_t / mean_c
            log_ratio = np.log(ratio)

            # Delta method variance
            var_log = var_t / (mean_t ** 2) + var_c / (mean_c ** 2)
            se_log = np.sqrt(var_log)

            ci_log_lower = log_ratio - z_crit * se_log
            ci_log_upper = log_ratio + z_crit * se_log

            # Back-transform
            ci_lower = np.exp(ci_log_lower)
            ci_upper = np.exp(ci_log_upper)

            return {
                'transformation': 'log',
                'ratio': ratio,
                'log_ratio': log_ratio,
                'confidence_interval': {
                    'level': confidence_level,
                    'lower': ci_lower,
                    'upper': ci_upper
                }
            }

        elif transformation == 'logit':
            # Logit transformation (for proportions)
            # p must be in (0, 1)
            def logit(p):
                return np.log(p / (1 - p))

            def inv_logit(x):
                return 1 / (1 + np.exp(-x))

            logit_t = logit(mean_t)
            logit_c = logit(mean_c)
            diff_logit = logit_t - logit_c

            # Delta method variance
            var_logit = var_t / (mean_t ** 2 * (1 - mean_t) ** 2) + \
                       var_c / (mean_c ** 2 * (1 - mean_c) ** 2)
            se_logit = np.sqrt(var_logit)

            ci_logit_lower = diff_logit - z_crit * se_logit
            ci_logit_upper = diff_logit + z_crit * se_logit

            # Back-transform
            ci_lower = inv_logit(ci_logit_lower)
            ci_upper = inv_logit(ci_logit_upper)

            return {
                'transformation': 'logit',
                'difference_logit': diff_logit,
                'confidence_interval': {
                    'level': confidence_level,
                    'lower': ci_lower,
                    'upper': ci_upper
                }
            }

        else:
            raise ValueError(f"Unknown transformation: {transformation}")

\end{lstlisting}

\section{Result Interpretation}

Statistical results must be translated into business insights and actionable recommendations.

\subsection{Statistical vs. Practical Significance}

\textbf{Statistical significance} indicates whether an effect is likely real (not due to chance). \textbf{Practical significance} indicates whether the effect is large enough to matter for business decisions.

\begin{example}[Statistical vs. Practical Significance]
An A/B test with 1 million users finds that variant B increases conversion rate from 10.0\% to 10.05\% (p < 0.001).

\textbf{Statistical significance:} Yes (p < 0.001)

\textbf{Practical significance:} Depends on context
\begin{itemize}
    \item If incremental implementation cost is \$10,000 and each conversion worth \$50, then: Additional revenue = $1,000,000 \times 0.0005 \times \$50 = \$25,000$ per million users
    \item ROI = $(\$25,000 - \$10,000) / \$10,000 = 150\%$ $\rightarrow$ \textbf{Practically significant}
\end{itemize}

\textbf{Decision:} Ship despite small effect size due to favorable cost-benefit.
\end{example}

\lstset{language=Python, caption=Business Impact Assessment}
\begin{lstlisting}
class BusinessImpactAnalyzer:
    """Assess business impact of A/B test results."""

    @staticmethod
    def calculate_business_impact(test_results: Dict,
                                 annual_users: int,
                                 value_per_conversion: float,
                                 implementation_cost: float) -> Dict:
        """
        Calculate business impact metrics.

        Parameters:
        -----------
        test_results : Dict
            Results from statistical test (must include difference and CI)
        annual_users : int
            Expected annual users
        value_per_conversion : float
            Dollar value per conversion/success
        implementation_cost : float
            One-time implementation cost

        Returns:
        --------
        Dict with business impact assessment
        """
        # Extract effect size
        if 'absolute_difference' in test_results:
            effect = test_results['absolute_difference']
        elif 'difference' in test_results:
            effect = test_results['difference']
        else:
            raise ValueError("Test results must include effect size")

        ci = test_results.get('confidence_interval', {})
        ci_lower = ci.get('lower', None)
        ci_upper = ci.get('upper', None)

        # Annual impact
        annual_additional_conversions = annual_users * effect
        annual_revenue_impact = annual_additional_conversions * value_per_conversion

        # ROI calculation
        roi = (annual_revenue_impact - implementation_cost) / implementation_cost * 100 if implementation_cost > 0 else np.inf

        # Conservative estimate (lower CI bound)
        if ci_lower is not None:
            conservative_conversions = annual_users * ci_lower
            conservative_revenue = conservative_conversions * value_per_conversion
            conservative_roi = (conservative_revenue - implementation_cost) / implementation_cost * 100 if implementation_cost > 0 else np.inf
        else:
            conservative_revenue = None
            conservative_roi = None

        # Optimistic estimate (upper CI bound)
        if ci_upper is not None:
            optimistic_conversions = annual_users * ci_upper
            optimistic_revenue = optimistic_conversions * value_per_conversion
            optimistic_roi = (optimistic_revenue - implementation_cost) / implementation_cost * 100 if implementation_cost > 0 else np.inf
        else:
            optimistic_revenue = None
            optimistic_roi = None

        # Practical significance assessment
        min_meaningful_effect = implementation_cost / (annual_users * value_per_conversion)
        is_practically_significant = abs(effect) > min_meaningful_effect

        return {
            'effect_size': effect,
            'annual_additional_conversions': annual_additional_conversions,
            'expected_annual_revenue_impact': annual_revenue_impact,
            'implementation_cost': implementation_cost,
            'expected_roi_pct': roi,
            'conservative_estimate': {
                'annual_revenue_impact': conservative_revenue,
                'roi_pct': conservative_roi
            },
            'optimistic_estimate': {
                'annual_revenue_impact': optimistic_revenue,
                'roi_pct': optimistic_roi
            },
            'min_meaningful_effect': min_meaningful_effect,
            'is_practically_significant': is_practically_significant,
            'is_statistically_significant': test_results.get('significant', False)
        }

    @staticmethod
    def create_impact_summary(impact_results: Dict) -> str:
        """Generate human-readable impact summary."""
        summary = []
        summary.append("=" * 60)
        summary.append("BUSINESS IMPACT SUMMARY")
        summary.append("=" * 60)
        summary.append("")

        summary.append(f"Effect Size: {impact_results['effect_size']:.4f}")
        summary.append(f"Annual Additional Conversions: {impact_results['annual_additional_conversions']:,.0f}")
        summary.append("")

        summary.append(f"Expected Annual Revenue Impact: ${impact_results['expected_annual_revenue_impact']:,.2f}")
        summary.append(f"Implementation Cost: ${impact_results['implementation_cost']:,.2f}")
        summary.append(f"Expected ROI: {impact_results['expected_roi_pct']:.1f}%")
        summary.append("")

        if impact_results['conservative_estimate']['annual_revenue_impact']:
            summary.append(f"Conservative Estimate (95% CI lower bound):")
            summary.append(f"  Revenue Impact: ${impact_results['conservative_estimate']['annual_revenue_impact']:,.2f}")
            summary.append(f"  ROI: {impact_results['conservative_estimate']['roi_pct']:.1f}%")
            summary.append("")

        summary.append(f"Statistical Significance: {'Yes' if impact_results['is_statistically_significant'] else 'No'}")
        summary.append(f"Practical Significance: {'Yes' if impact_results['is_practically_significant'] else 'No'}")
        summary.append("")

        # Decision recommendation
        if impact_results['is_statistically_significant'] and impact_results['is_practically_significant']:
            recommendation = "RECOMMEND: Ship - Both statistically and practically significant"
        elif impact_results['is_statistically_significant'] and not impact_results['is_practically_significant']:
            recommendation = "CONSIDER: Small effect - Evaluate strategic value beyond ROI"
        elif not impact_results['is_statistically_significant'] and impact_results['is_practically_significant']:
            recommendation = "INCONCLUSIVE: Run longer test for more power"
        else:
            recommendation = "DO NOT SHIP: Not significant"

        summary.append(f"Recommendation: {recommendation}")
        summary.append("")

        return "\n".join(summary)

\end{lstlisting}

\subsection{Segment Analysis}

Treatment effects may vary across user segments. Segmented analysis reveals heterogeneous treatment effects (HTEs).

\lstset{language=Python, caption=Segment Analysis Framework}
\begin{lstlisting}
class SegmentAnalyzer:
    """Analyze treatment effects across segments."""

    @staticmethod
    def segment_analysis(data: pd.DataFrame,
                        metric_col: str,
                        segment_col: str,
                        variant_col: str = 'variant',
                        test_type: str = 'auto') -> Dict:
        """
        Analyze treatment effect within each segment.

        Parameters:
        -----------
        data : DataFrame
            Experimental data
        metric_col : str
            Metric to analyze
        segment_col : str
            Column defining segments
        variant_col : str
            Variant column
        test_type : str
            'auto', 't_test', 'mann_whitney', or 'proportion'

        Returns:
        --------
        Dict with segment-specific results
        """
        executor = ABTestExecutor()
        segment_results = {}

        for segment in data[segment_col].unique():
            segment_data = data[data[segment_col] == segment]

            # Run appropriate test
            if test_type == 'auto':
                result = executor.automatic_test_selection(
                    segment_data, metric_col, variant_col
                )
            elif test_type == 't_test':
                result = executor.t_test(segment_data, metric_col, variant_col)
            elif test_type == 'mann_whitney':
                result = executor.mann_whitney_test(segment_data, metric_col, variant_col)
            elif test_type == 'proportion':
                result = executor.proportion_test(segment_data, metric_col, variant_col)
            else:
                raise ValueError(f"Unknown test_type: {test_type}")

            segment_results[segment] = result

        return segment_results

    @staticmethod
    def test_for_interaction(data: pd.DataFrame,
                            metric_col: str,
                            segment_col: str,
                            variant_col: str = 'variant') -> Dict:
        """
        Test if treatment effect differs significantly across segments.

        Uses two-way ANOVA or similar to test for interaction.

        Returns:
        --------
        Dict with interaction test results
        """
        from scipy import stats as sp_stats
        import statsmodels.formula.api as smf

        # Create interaction term
        data_test = data.copy()
        data_test['interaction'] = (
            data_test[variant_col].astype(str) + '_' +
            data_test[segment_col].astype(str)
        )

        # Two-way ANOVA
        formula = f'{metric_col} ~ C({variant_col}) + C({segment_col}) + C({variant_col}):C({segment_col})'

        try:
            model = smf.ols(formula, data=data_test).fit()
            anova_table = sm.stats.anova_lm(model, typ=2)

            interaction_p = anova_table.loc[
                f'C({variant_col}):C({segment_col})', 'PR(>F)'
            ]

            return {
                'has_interaction': interaction_p < 0.05,
                'interaction_p_value': interaction_p,
                'anova_table': anova_table,
                'interpretation': (
                    "Treatment effect varies significantly across segments"
                    if interaction_p < 0.05
                    else "No significant interaction detected"
                )
            }
        except Exception as e:
            return {
                'error': str(e),
                'message': "Could not perform interaction test"
            }

    @staticmethod
    def visualize_segments(data: pd.DataFrame,
                          metric_col: str,
                          segment_col: str,
                          variant_col: str = 'variant'):
        """Create segment analysis visualization."""
        import matplotlib.pyplot as plt
        import seaborn as sns

        segments = data[segment_col].unique()
        n_segments = len(segments)

        fig, axes = plt.subplots(1, min(n_segments, 4), figsize=(16, 4))
        if n_segments == 1:
            axes = [axes]

        for idx, segment in enumerate(segments[:4]):
            segment_data = data[data[segment_col] == segment]

            # Create violin + box plot
            sns.violinplot(
                data=segment_data,
                x=variant_col,
                y=metric_col,
                ax=axes[idx],
                alpha=0.5
            )

            # Calculate means and plot
            segment_stats = DescriptiveStats()
            ci_data = segment_stats.confidence_intervals(
                segment_data, metric_col, variant_col
            )

            for i, row in ci_data.iterrows():
                axes[idx].plot(
                    [i, i],
                    [row['ci_lower'], row['ci_upper']],
                    'k-',
                    linewidth=2
                )
                axes[idx].plot(i, row['mean'], 'ro', markersize=8)

            axes[idx].set_title(f'Segment: {segment}')
            axes[idx].set_ylabel(metric_col)
            axes[idx].grid(alpha=0.3)

        plt.tight_layout()
        return fig

    @staticmethod
    def create_segment_report(segment_results: Dict) -> pd.DataFrame:
        """Create summary table of segment results."""
        rows = []

        for segment, result in segment_results.items():
            if 'control_mean' in result:
                control_val = result['control_mean']
                treatment_val = result['treatment_mean']
            elif 'control_rate' in result:
                control_val = result['control_rate']
                treatment_val = result['treatment_rate']
            elif 'control_median' in result:
                control_val = result['control_median']
                treatment_val = result['treatment_median']
            else:
                continue

            rows.append({
                'segment': segment,
                'control': control_val,
                'treatment': treatment_val,
                'difference': result.get('difference',  result.get('absolute_difference', treatment_val - control_val)),
                'p_value': result.get('p_value', np.nan),
                'significant': result.get('significant', False)
            })

        return pd.DataFrame(rows)

\end{lstlisting}

\subsection{Long-term Effects Consideration}

Short-term test results may not reflect long-term impact.

\begin{itemize}
    \item \textbf{Novelty effects:} Users may engage more/less with new features initially
    \item \textbf{Learning effects:} Behavior changes as users become familiar
    \item \textbf{Network effects:} Impact may compound as more users adopt
    \item \textbf{Cannibalization:} Short-term gains may reduce long-term metrics
\end{itemize}

\textbf{Strategies for assessing long-term effects:}
\begin{enumerate}
    \item Run experiments for multiple weeks/months
    \item Analyze cohorts based on exposure duration
    \item Monitor post-experiment metrics
    \item Use holdout groups for long-term tracking
\end{enumerate}

\section{Reporting and Communication}

Effective communication ensures stakeholders understand results and can make informed decisions.

\subsection{Stakeholder Communication Strategies}

Different audiences require different levels of detail:

\begin{center}
\begin{tabular}{lll}
\toprule
\textbf{Audience} & \textbf{Focus} & \textbf{Format} \\
\midrule
Executives & Business impact, ROI & 1-page summary \\
Product managers & Metrics, segments, recommendations & Dashboard + slides \\
Engineers & Implementation details, edge cases & Technical doc \\
Data scientists & Methods, assumptions, limitations & Full analysis \\
\bottomrule
\end{tabular}
\end{center}

\subsection{Visualization Best Practices}

\lstset{language=Python, caption=Professional Result Visualization}
\begin{lstlisting}
class ResultVisualizer:
    """Create professional visualizations for reporting."""

    @staticmethod
    def create_executive_summary_viz(test_results: Dict,
                                     metric_name: str):
        """Create executive summary visualization."""
        import matplotlib.pyplot as plt

        fig, axes = plt.subplots(1, 2, figsize=(14, 5))

        # 1. Treatment effect with CI
        if 'control_mean' in test_results:
            control_val = test_results['control_mean']
            treatment_val = test_results['treatment_mean']
            value_type = 'Mean'
        elif 'control_rate' in test_results:
            control_val = test_results['control_rate']
            treatment_val = test_results['treatment_rate']
            value_type = 'Rate'
        else:
            raise ValueError("Unknown result format")

        ci = test_results['confidence_interval']

        # Bar chart with error bars
        variants = ['Control', 'Treatment']
        values = [control_val, treatment_val]

        bars = axes[0].bar(variants, values, color=['#3498db', '#e74c3c'], alpha=0.7)

        # Add CI for treatment
        diff = test_results.get('difference', test_results.get('absolute_difference'))
        axes[0].errorbar(
            [1],
            [treatment_val],
            yerr=[[treatment_val - ci['lower']], [ci['upper'] - treatment_val]],
            fmt='none',
            color='black',
            capsize=10,
            linewidth=2
        )

        # Annotations
        for i, (bar, val) in enumerate(zip(bars, values)):
            height = bar.get_height()
            axes[0].text(
                bar.get_x() + bar.get_width()/2.,
                height,
                f'{val:.4f}',
                ha='center',
                va='bottom',
                fontsize=12,
                fontweight='bold'
            )

        axes[0].set_ylabel(f'{value_type} {metric_name}', fontsize=12)
        axes[0].set_title(f'{metric_name} Comparison', fontsize=14, fontweight='bold')
        axes[0].grid(axis='y', alpha=0.3)

        # 2. Effect size visualization
        rel_diff = test_results.get('relative_difference_pct', 0)

        color = '#27ae60' if rel_diff > 0 else '#e74c3c'

        axes[1].barh(['Relative\nLift'], [rel_diff], color=color, alpha=0.7)
        axes[1].axvline(0, color='black', linestyle='--', linewidth=1)

        axes[1].text(
            rel_diff,
            0,
            f'{rel_diff:+.2f}%',
            ha='left' if rel_diff > 0 else 'right',
            va='center',
            fontsize=14,
            fontweight='bold'
        )

        axes[1].set_xlabel('Percentage Change', fontsize=12)
        axes[1].set_title('Treatment Effect Size', fontsize=14, fontweight='bold')
        axes[1].grid(axis='x', alpha=0.3)

        # Overall title
        sig_status = "Significant" if test_results.get('significant') else "Not Significant"
        p_val = test_results.get('p_value', np.nan)

        fig.suptitle(
            f'A/B Test Results: {metric_name}\n{sig_status} (p={p_val:.4f})',
            fontsize=16,
            fontweight='bold'
        )

        plt.tight_layout()
        return fig

    @staticmethod
    def create_comprehensive_report_viz(data: pd.DataFrame,
                                       metric_col: str,
                                       test_results: Dict,
                                       variant_col: str = 'variant'):
        """Create comprehensive report with multiple visualizations."""
        import matplotlib.pyplot as plt
        import seaborn as sns

        fig = plt.figure(figsize=(16, 12))
        gs = fig.add_gridspec(3, 2, hspace=0.3, wspace=0.3)

        # 1. Distribution comparison
        ax1 = fig.add_subplot(gs[0, :])
        for variant in data[variant_col].unique():
            variant_data = data[data[variant_col] == variant][metric_col].dropna()
            sns.kdeplot(data=variant_data, label=str(variant), ax=ax1, linewidth=2)
        ax1.set_xlabel(metric_col, fontsize=11)
        ax1.set_ylabel('Density', fontsize=11)
        ax1.set_title('Distribution Comparison', fontsize=13, fontweight='bold')
        ax1.legend()
        ax1.grid(alpha=0.3)

        # 2. Effect size with CI
        ax2 = fig.add_subplot(gs[1, 0])

        diff = test_results.get('difference', test_results.get('absolute_difference', 0))
        ci = test_results.get('confidence_interval', {})

        ax2.plot([0], [diff], 'o', markersize=12, color='#e74c3c')
        if ci:
            ax2.plot(
                [0, 0],
                [ci.get('lower', diff), ci.get('upper', diff)],
                linewidth=3,
                color='#3498db'
            )

        ax2.axhline(0, color='gray', linestyle='--', linewidth=1)
        ax2.set_xlim(-0.5, 0.5)
        ax2.set_ylabel(f'Treatment Effect on {metric_col}', fontsize=11)
        ax2.set_title('Effect Size with 95% CI', fontsize=13, fontweight='bold')
        ax2.set_xticks([])
        ax2.grid(axis='y', alpha=0.3)

        # 3. Sample sizes
        ax3 = fig.add_subplot(gs[1, 1])

        sample_sizes = test_results.get('sample_sizes', {})
        if sample_sizes:
            variants_list = list(sample_sizes.keys())
            sizes = list(sample_sizes.values())

            bars = ax3.bar(variants_list, sizes, color=['#3498db', '#e74c3c'], alpha=0.7)

            for bar, size in zip(bars, sizes):
                height = bar.get_height()
                ax3.text(
                    bar.get_x() + bar.get_width()/2.,
                    height,
                    f'n={size:,}',
                    ha='center',
                    va='bottom',
                    fontsize=10
                )

            ax3.set_ylabel('Sample Size', fontsize=11)
            ax3.set_title('Sample Sizes by Variant', fontsize=13, fontweight='bold')
            ax3.grid(axis='y', alpha=0.3)

        # 4. Statistical summary table
        ax4 = fig.add_subplot(gs[2, :])
        ax4.axis('tight')
        ax4.axis('off')

        summary_data = [
            ['Metric', metric_col],
            ['Test Type', test_results.get('test_type', 'N/A')],
            ['P-value', f"{test_results.get('p_value', np.nan):.4f}"],
            ['Significance Level', '0.05'],
            ['Significant', 'Yes' if test_results.get('significant') else 'No'],
            ['Effect Size', f"{diff:.4f}"],
            ['Relative Change', f"{test_results.get('relative_difference_pct', 0):.2f}%"],
            ['CI Lower', f"{ci.get('lower', np.nan):.4f}"],
            ['CI Upper', f"{ci.get('upper', np.nan):.4f}"]
        ]

        table = ax4.table(
            cellText=summary_data,
            cellLoc='left',
            loc='center',
            colWidths=[0.4, 0.6]
        )
        table.auto_set_font_size(False)
        table.set_fontsize(10)
        table.scale(1, 2)
        ax4.set_title('Statistical Summary', fontsize=13, fontweight='bold', pad=20)

        plt.suptitle(
            f'Comprehensive A/B Test Report: {metric_col}',
            fontsize=16,
            fontweight='bold',
            y=0.995
        )

        return fig

\end{lstlisting}

\subsection{Executive Summary Template}

\lstset{language=Python, caption=Executive Summary Generation}
\begin{lstlisting}
class ReportGenerator:
    """Generate professional A/B test reports."""

    @staticmethod
    def generate_executive_summary(test_results: Dict,
                                   metric_name: str,
                                   experiment_name: str,
                                   business_impact: Dict = None) -> str:
        """
        Generate one-page executive summary.

        Parameters:
        -----------
        test_results : Dict
            Statistical test results
        metric_name : str
            Name of primary metric
        experiment_name : str
            Name/description of experiment
        business_impact : Dict
            Business impact assessment results

        Returns:
        --------
        str : Formatted executive summary
        """
        summary = []
        summary.append("=" * 70)
        summary.append(f"EXECUTIVE SUMMARY: {experiment_name}")
        summary.append("=" * 70)
        summary.append("")

        # Key findings
        summary.append("KEY FINDINGS")
        summary.append("-" * 70)

        sig_status = "SIGNIFICANT" if test_results.get('significant') else "NOT SIGNIFICANT"
        direction = "increased" if test_results.get('difference', 0) > 0 else "decreased"

        rel_change = test_results.get('relative_difference_pct', 0)

        summary.append(
            f"The treatment {direction} {metric_name} by {abs(rel_change):.2f}% "
            f"(p={test_results.get('p_value', np.nan):.4f})"
        )
        summary.append(f"Result is {sig_status} at 95% confidence level")
        summary.append("")

        # Business impact
        if business_impact:
            summary.append("BUSINESS IMPACT")
            summary.append("-" * 70)

            revenue = business_impact.get('expected_annual_revenue_impact', 0)
            roi = business_impact.get('expected_roi_pct', 0)

            summary.append(f"Expected Annual Revenue Impact: ${revenue:,.0f}")
            summary.append(f"Expected ROI: {roi:.1f}%")

            if business_impact.get('conservative_estimate', {}).get('annual_revenue_impact'):
                cons_revenue = business_impact['conservative_estimate']['annual_revenue_impact']
                summary.append(f"Conservative Estimate (95% CI lower): ${cons_revenue:,.0f}")

            summary.append("")

        # Recommendation
        summary.append("RECOMMENDATION")
        summary.append("-" * 70)

        if business_impact:
            if (business_impact.get('is_statistically_significant') and
                business_impact.get('is_practically_significant')):
                recommendation = "SHIP: Both statistically and practically significant"
            elif (business_impact.get('is_statistically_significant') and
                  not business_impact.get('is_practically_significant')):
                recommendation = "CONSIDER: Statistically significant but small business impact"
            elif (not business_impact.get('is_statistically_significant') and
                  business_impact.get('is_practically_significant')):
                recommendation = "EXTEND TEST: Promising but needs more data"
            else:
                recommendation = "DO NOT SHIP: Not significant"
        else:
            recommendation = f"SHIP if {sig_status}" if test_results.get('significant') else "DO NOT SHIP"

        summary.append(recommendation)
        summary.append("")

        # Technical details
        summary.append("TECHNICAL DETAILS")
        summary.append("-" * 70)

        sample_sizes = test_results.get('sample_sizes', {})
        if sample_sizes:
            for variant, size in sample_sizes.items():
                summary.append(f"  {variant}: n={size:,}")

        summary.append(f"  Test type: {test_results.get('test_type', test_results.get('recommendation', 'N/A'))}")
        summary.append("")

        return "\n".join(summary)

    @staticmethod
    def generate_technical_documentation(test_results: Dict,
                                        data_summary: Dict,
                                        assumptions: Dict,
                                        metric_name: str) -> str:
        """Generate detailed technical documentation."""
        doc = []
        doc.append("=" * 70)
        doc.append("TECHNICAL DOCUMENTATION: A/B TEST ANALYSIS")
        doc.append("=" * 70)
        doc.append("")

        # 1. Data Summary
        doc.append("1. DATA SUMMARY")
        doc.append("-" * 70)
        doc.append(f"Primary Metric: {metric_name}")
        doc.append(f"Total Observations: {data_summary.get('n_total', 'N/A'):,}")
        doc.append(f"Date Range: {data_summary.get('date_range', 'N/A')}")
        doc.append("")

        # 2. Statistical Method
        doc.append("2. STATISTICAL METHOD")
        doc.append("-" * 70)
        doc.append(f"Test Type: {test_results.get('test_type', 'N/A')}")
        doc.append(f"Significance Level: α = 0.05")
        doc.append(f"Hypothesis: Two-sided")
        doc.append("")

        # 3. Assumptions
        doc.append("3. ASSUMPTIONS")
        doc.append("-" * 70)

        if assumptions:
            if 'normality' in assumptions:
                doc.append("Normality:")
                for variant, result in assumptions['normality'].items():
                    is_normal = result.get('shapiro_wilk', {}).get('is_normal', 'Unknown')
                    doc.append(f"  {variant}: {'Normal' if is_normal else 'Non-normal'}")

            if 'variance_equality' in assumptions:
                equal_var = assumptions['variance_equality'].get('levene_test', {}).get('equal_variances', 'Unknown')
                doc.append(f"Equal Variances: {'Yes' if equal_var else 'No'}")

        doc.append("")

        # 4. Results
        doc.append("4. RESULTS")
        doc.append("-" * 70)

        p_val = test_results.get('p_value', np.nan)
        doc.append(f"P-value: {p_val:.6f}")

        if 'control_mean' in test_results:
            doc.append(f"Control Mean: {test_results['control_mean']:.4f}")
            doc.append(f"Treatment Mean: {test_results['treatment_mean']:.4f}")
        elif 'control_rate' in test_results:
            doc.append(f"Control Rate: {test_results['control_rate']:.4f}")
            doc.append(f"Treatment Rate: {test_results['treatment_rate']:.4f}")

        diff = test_results.get('difference', test_results.get('absolute_difference', 0))
        doc.append(f"Absolute Difference: {diff:.4f}")
        doc.append(f"Relative Difference: {test_results.get('relative_difference_pct', 0):.2f}%")

        ci = test_results.get('confidence_interval', {})
        if ci:
            doc.append(f"95% CI: [{ci.get('lower', np.nan):.4f}, {ci.get('upper', np.nan):.4f}]")

        doc.append("")

        # 5. Interpretation
        doc.append("5. INTERPRETATION")
        doc.append("-" * 70)

        if test_results.get('significant'):
            doc.append("The treatment has a statistically significant effect.")
        else:
            doc.append("No statistically significant difference detected.")

        doc.append("")

        return "\n".join(doc)

\end{lstlisting}

\section{Summary}

This chapter presented a comprehensive framework for analyzing and interpreting A/B test results:

\begin{itemize}
    \item \textbf{Data Preprocessing:} Systematic validation, outlier treatment, missing data handling, and transformations ensure data quality
    \item \textbf{Exploratory Analysis:} Descriptive statistics and visualizations reveal patterns before formal testing
    \item \textbf{Statistical Pipeline:} Automated assumption checking and test selection ensure appropriate methodology
    \item \textbf{Effect Sizes:} Cohen's d, relative effects, and other measures quantify practical significance
    \item \textbf{Business Impact:} ROI calculations connect statistical findings to business decisions
    \item \textbf{Segment Analysis:} Heterogeneous treatment effects reveal differential impacts
    \item \textbf{Communication:} Tailored reporting ensures all stakeholders can act on findings
\end{itemize}

The Python frameworks provided enable reproducible, production-ready analysis workflows. By combining statistical rigor with business context, we transform experimental data into actionable insights.

\section{Further Reading}

\begin{itemize}
    \item \cite{deng2016data}: Data quality in online experiments
    \item \cite{kohavi2020trustworthy}: Comprehensive experimentation guide
    \item \cite{gelman2006data}: Data analysis using regression and multilevel models
    \item \cite{harrell2015regression}: Regression modeling strategies
    \item \cite{tukey1977exploratory}: Exploratory data analysis foundations
\end{itemize}

\section{Exercises}

\begin{enumerate}
    \item \textbf{Data Preprocessing Pipeline:} Implement a complete preprocessing pipeline for A/B test data including:
    \begin{itemize}
        \item SRM detection
        \item Outlier identification using multiple methods
        \item Missing data imputation
        \item Transformation assessment
    \end{itemize}

    \item \textbf{EDA Dashboard:} Create an interactive EDA dashboard using Plotly that includes:
    \begin{itemize}
        \item Distribution comparisons
        \item Temporal trends
        \item Statistical summaries
        \item Assumption diagnostics
    \end{itemize}

    \item \textbf{Automated Analysis:} Build an automated analysis pipeline that:
    \begin{itemize}
        \item Selects appropriate statistical tests based on data characteristics
        \item Calculates effect sizes and confidence intervals
        \item Assesses business impact
        \item Generates executive summary
    \end{itemize}

    \item \textbf{Segment Analysis:} Analyze treatment effects across multiple segments and:
    \begin{itemize}
        \item Test for interactions
        \item Visualize heterogeneous effects
        \item Report segment-specific recommendations
    \end{itemize}

    \item \textbf{Report Generation:} Create templates for:
    \begin{itemize}
        \item Executive one-pagers
        \item Technical documentation
        \item Stakeholder presentations
    \end{itemize}
\end{enumerate}
